\documentclass[12pt,UTF8,fontset=none,twocolumn]{ctexart}
% 使用「繁體中文」常見 Windows 字型，避免 Fandol 缺字（如「戶、為、線」等）。
% 若本機缺少字型，會自動退回到 Noto/Source Han 或 Fandol。
%
% 優先使用微軟正黑體（字形完整、含粗體），其次使用新細明體/細明體。
\IfFontExistsTF{Microsoft JhengHei}{
  \setCJKmainfont{Microsoft JhengHei}
  \setCJKsansfont{Microsoft JhengHei}
}{
  \IfFontExistsTF{PMingLiU}{
    \setCJKmainfont{PMingLiU}
    \IfFontExistsTF{Microsoft JhengHei}{\setCJKsansfont{Microsoft JhengHei}}{}
  }{
    \IfFontExistsTF{MingLiU}{
      \setCJKmainfont{MingLiU}
      \IfFontExistsTF{Microsoft JhengHei}{\setCJKsansfont{Microsoft JhengHei}}{}
    }{
      \IfFontExistsTF{Noto Serif CJK TC}{
        \setCJKmainfont{Noto Serif CJK TC}
      }{
        \IfFontExistsTF{Source Han Serif TC}{
          \setCJKmainfont{Source Han Serif TC}
        }{
          \setCJKmainfont{FandolSong-Regular}
        }
      }
      \IfFontExistsTF{Noto Sans CJK TC}{
        \setCJKsansfont{Noto Sans CJK TC}
      }{
        \IfFontExistsTF{Source Han Sans TC}{
          \setCJKsansfont{Source Han Sans TC}
        }{
          \setCJKsansfont{FandolHei-Regular}
        }
      }
    }
  }
}

\usepackage{geometry}
\usepackage{microtype}
\usepackage{amsmath,amssymb}
\usepackage{booktabs}
\usepackage{graphicx}
\usepackage{enumitem}
\usepackage{hyperref}
\usepackage{cite}
\usepackage{xcolor}
\usepackage{float}
\usepackage{algorithm}
\usepackage{algpseudocode}

% Algorithm styling (ruled + line numbers + blue // comments)
\floatstyle{ruled}
\restylefloat{algorithm}
\algrenewcommand\algorithmicrequire{\textbf{Input:}}
\algrenewcommand\algorithmicensure{\textbf{Output:}}
\algrenewcommand\algorithmiccomment[1]{\hfill{\color{blue}\ttfamily // #1}}

\geometry{margin=1in}
\setlength{\columnsep}{0.25in}
\hypersetup{
  colorlinks=true,
  linkcolor=blue,
  citecolor=blue,
  urlcolor=blue
}

\makeatletter
% Section numbering: I., II., ... and subsections: A., B., ...
\renewcommand{\thesection}{\Roman{section}}
\renewcommand{\thesubsection}{\Alph{subsection}}
\renewcommand{\@seccntformat}[1]{\csname the#1\endcsname.\quad}
\makeatother

\title{Seamless Coverage and ISL-Assisted Ferrying for Sustained User Satisfaction in NGSO Constellations}
\author{Chia-Yu Hu and Guey-Yun Ghang}
\date{}

\begin{document}
\maketitle
\thispagestyle{plain}
\pagestyle{plain}

\section{Introduction}
\label{sec:introduction}
近年 NGSO/LEO constellations 快速部署，使全球性寬頻覆蓋與低延遲通訊成為可能，並被視為未來整合式非地面網路（NTN）與 6G 架構的重要組成 \cite{giordani2021ntn,lin2021leosat6g,zhujiang2022istn}。此外，NGSO systems 的通訊架構、資源管理與與地面網路整合等議題亦已被系統性整理 \cite{alhraishawi2022ngsoSurvey}。在寬頻衛星系統常用的 Ku/Ka bands 中，NGSO 與既有 GSO networks 往往需在相同或相鄰頻段共享資源（例如 fixed-satellite service, FSS），因此 co-channel interference 成為不可避免的關鍵限制。當 constellation 規模擴大到數千顆衛星時，多顆衛星的 \emph{aggregated emissions} 可能對受保護的 GSO earth stations 造成顯著干擾，且干擾強度與衛星幾何位置高度相關並隨時間快速變動。

為保護既有 GSO networks，International Telecommunication Union（ITU）在 \emph{Radio Regulations} 的 Article~22 制定 NGSO--GSO 共存下的干擾限制，並以 equivalent power flux-density（EPFD）等指標評估 NGSO constellation 對 GSO earth stations 的\emph{聚合}干擾 \cite{itur_rr2020}. 由於 EPFD 規範需要可重現且一致的干擾計算流程，ITU-R 亦提供用於符合性檢核的功能性描述（S.1503-3），來判定 NGSO/FSS 系統是否符合 Article~22 的限制 \cite{itur_s1503_2018}. 在實際干擾評估與資源最佳化建模中，天線主/旁瓣與離軸角（off-axis angle）對干擾有決定性影響，因此常採用 ITU-R 所建議的輻射方向圖模型：例如 S.1528 用於 30~GHz 以下的 NGSO 衛星天線方向圖 \cite{itur_s1528_2001}，以及 S.1428-1 作為 10.7--30~GHz 頻段非 GSO 干擾評估所用的 FSS earth station 參考方向圖 \cite{itur_s1428_2001}。上述法規與建模框架共同形成 NGSO constellation 在 Ka-band 等頻段進行干擾緩解與服務維持時必須遵循的基本約束。

因此，NGSO resource management 本質上具有多目標特性：一方面必須符合 EPFD regulatory constraints，另一方面又要最大化用戶服務能力（例如 throughput、coverage 或 demand satisfaction）。既有研究指出，在 dense NGSO deployments 下的 interference risk 需要被量化評估，且可能涉及 operator coordination 與 mitigation trade-offs \cite{braun2019dyspan}。在實務上，典型 mitigation 手段包含 power control、tilt/pitch 調整與空間隔離等\cite{hills2022beamsteering}。在 EPFD 約束下，Jalali 等人進一步提出在 constellation 尺度聯合最佳化 transmit power 與 tilt/pitch，以提升 demand satisfaction 並降低干擾風險 \cite{jalali2023joint}。另一方面，Li 等人針對 OneWeb 的 progressive pitch 建立模型並反推可行參數，指出該策略可在多數地區有效降低干擾，但仍可能在部分時間產生 service blind spots \cite{li2021progressive}。這些工作凸顯了 interference mitigation 與 service continuity 之間的權衡：當衛星為避免干擾而調整 pitch 或關閉部分 beams 時，其有效 coverage footprint 可能縮小，導致使用者短暫失去覆蓋，進而降低整體 satisfaction 與使用體驗。

本論文聚焦於在干擾約束下維持 \emph{seamless service}。相關研究已指出，對採用 progressive pitch 的 LEO constellations，容易在相鄰衛星覆蓋銜接處產生 blind spots。我們將干擾事件中幾何條件最敏感、最可能主導 EPFD 違規的衛星稱為 \emph{critical satellites}，並探討當 critical satellites 必須調整發射行為時，如何避免造成使用者 zero-coverage 與 satisfaction 下降。除了下行發射端的 power/tilt 控制外，現代巨型 constellations 亦逐步導入 inter-satellite links (ISL)，可在星座內部提供額外的 routing 與 load balancing 自由度。本文據此提出 \emph{seamless coverage} 與 \emph{ISL-assisted ferrying} 兩項設計，降低 interference mitigation 對 service continuity 的負面影響。

\subsection{Contributions}
\label{subsec:contributions}
本論文主要貢獻如下：
\begin{itemize}[leftmargin=*,itemsep=0.25em]
  \item \textbf{Seamless coverage constraint.} 提出 seamless coverage 設計目標，確保每一位使用者在任意時間至少被一顆衛星覆蓋，避免因 interference mitigation（如 tilt/pitch 調整或 power reduction）導致的 zero-coverage intervals。
  \item \textbf{ISL-assisted ferrying via same-orbit north/south neighbors.} 針對 critical satellites，選擇同軌道上幾何距離最近的 north/south neighbors，透過 ISL 進行 traffic relaying，以提升 service continuity 並維持較高 demand satisfaction。
\end{itemize}

\section{Related Work}
\label{sec:related-work}
\subsection{Interference mitigation for NGSO--GSO coexistence}
NGSO--GSO coexistence 下的干擾問題已引發大量 interference analysis 與 mitigation 研究。常見的 spatial isolation 方法包含 arc avoidance、exclusion zones，以及 beam switching-off 等作法；這些方法雖能降低干擾，卻可能以犧牲 coverage 與 service continuity 為代價。較具彈性的作法則在干擾約束下進行 resource optimization。Jalali 等人以 ITU 定義的 EPFD 作為 aggregate interference limit，並聯合最佳化 constellation 的 transmit power 與 tilt/pitch，以在 regulatory compliance 的同時最大化 demand satisfaction \cite{jalali2023joint}。該工作指出 EPFD events 具有顯著的地理相依性與幾何敏感性，並可透過僅對少數 critical satellites 進行 tilt 以降低操作複雜度。然而，tilt 或 beam switching-off 仍可能縮小 coverage footprint，進而在某些時段造成 coverage gaps。

\subsection{OneWeb progressive pitch and service blind spots}
OneWeb 提出 progressive pitch 策略以保護 GSO networks，但其具體參數與細節並未完整公開。Li 等人針對 OneWeb--GSO coexistence 建立數學模型，在滿足 GSO interference protection criteria 下求解 progressive pitch parameters，並以數值模擬驗證其 interference avoidance 效果 \cite{li2021progressive}。該研究指出 progressive pitch 在全球多數區域具有良好 interference avoidance 能力，但在部分時間仍可能出現 service blind spots，並分析其成因與可能的消除方式。這說明：以 pitch adjustment / beam switching-off 為主的 mitigation 策略，仍可能在特定幾何條件下導致 service continuity 問題。

\subsection{Baselines}
本文以兩篇文獻作為主要 baselines：（i）在 aggregate EPFD constraints 下聯合最佳化 power 與 tilt/pitch 以提升 demand satisfaction 的方法 \cite{jalali2023joint}；（ii）針對 OneWeb progressive pitch 之參數反推與 interference avoidance 效果分析 \cite{li2021progressive}。本文之方法在此基礎上進一步引入 seamless coverage constraint 與 ISL-assisted ferrying，以降低 mitigation 造成的 service blind spots 並提升整體 satisfaction。

\subsection*{Research Gap}
現有 NGSO--GSO interference mitigation 多以滿足 EPFD constraints 為主，常透過 power/tilt 或 beam switching-off 達成；但這些做法可能改變 coverage footprint，造成 coverage gaps / service blind spots，降低 service continuity \cite{li2021progressive,jalali2023joint}。此外，EPFD 作為\emph{聚合}干擾指標具有幾何相依與時間變動特性，使得僅靠局部或靜態設計難以同時兼顧「合規」與「無縫覆蓋」\cite{itur_rr2020,itur_s1503_2018}。因此，本研究聚焦於在 EPFD 合規下顯式加入 seamless coverage，並利用 ISL 進行 relay/ferrying 以維持較高的 user satisfaction。

% \section{System Model and Problem Formulation}

\subsection{System Model}

\subsubsection{Network Scenario}

本研究考慮一個 OneWeb-like 的 LEO 衛星通訊系統，並聚焦於 NGSO--GSO coexistence 情境下的下行干擾管理問題。系統中每顆 LEO 衛星配置 $16$ 個 spot beams，並於每一個 time slot 中在衛星覆蓋範圍內隨機生成地面 users。設 $\mathcal{S}(t)$ 為當前 time slot $t$ 中可見之 LEO 衛星集合，$\mathcal{B}_s$ 為衛星 $s$ 的 beam 集合，且 $|\mathcal{B}_s|=16$。

在每一個 time slot 中，每位 user 僅能由一顆衛星的一個 beam 服務。設 $\mathcal{U}(t)$ 為當前時槽中的 user 集合，並以二元變數 $x_{u,s,b}(t)$ 表示 user--beam association。若 user $u$ 由衛星 $s$ 的 beam $b$ 服務，則 $x_{u,s,b}(t)=1$，否則為 $0$。由於每位 user 僅能對應到單一 serving satellite 與單一 beam，因此 association 必須滿足唯一性限制。

此外，系統中設有一個受保護之 GSO ground station (GS)。由於 LEO 衛星在特定時間可能與 GS 呈現近 inline 幾何關係，因此當某顆 LEO 衛星接近 GS 正上方時，其下行發射較容易造成較高的 aggregate EPFD。本文將該時刻對 GS 造成最大干擾風險的 LEO 衛星定義為 critical satellite，記為 $s^{\star}(t)$。同時，考慮 critical satellite 同軌道上的北側與南側鄰近衛星，分別記為 $s^{\mathrm{N}}(t)$ 與 $s^{\mathrm{S}}(t)$，作為後續 relay / reassociation 的候選 helper satellites。

\subsubsection{User Association and Service Model}

設由衛星 $s$ 的 beam $b$ 所服務之 user 集合為
\begin{equation}
\mathcal{U}_{s,b}(t)=\{u\in\mathcal{U}(t)\mid x_{u,s,b}(t)=1\}.
\end{equation}

則 beam 上的 user 數定義為
\begin{equation}
N_{s,b}(t)=\sum_{u\in\mathcal{U}(t)} x_{u,s,b}(t).
\end{equation}

每位 user 的需求速率記為 $D_u(t)$。若所有 users 具有相同需求，則可簡化為固定常數 $D$，例如本研究設定為 $25$ Mbps。

由於每位 user 僅能由單一衛星之單一 beam 服務，因此需滿足
\begin{equation}
\sum_{s\in\mathcal{S}(t)}\sum_{b\in\mathcal{B}_s} x_{u,s,b}(t)\le 1,
\qquad \forall u\in\mathcal{U}(t).
\end{equation}

此外，若 user $u$ 不在 beam $b$ 的可服務覆蓋區 $\mathcal{C}_{s,b}(t)$ 內，則不可建立 association，即
\begin{equation}
x_{u,s,b}(t)=0,
\qquad \forall u\notin\mathcal{C}_{s,b}(t).
\end{equation}

\subsubsection{Load-Aware Power Allocation}

每顆衛星 $s$ 具有固定最大總發射功率 $P_s^{\max}$。此外，考量實際多波束衛星酬載架構中，各 beam 的發射能力通常仍受限於 beam-level RF chain、power amplifier 或 EIRP 上限，因此本文進一步對每個 beam 設定最大可分配功率，記為 $P_{s,b}^{\max}$。

在初始 user association 完成後，本文依據 active beams 上的 user 數比例分配衛星總功率，以反映較高負載的 beam 應獲得較多傳輸功率。設 $\mathcal{B}_s^{\mathrm{act}}(t)$ 為衛星 $s$ 在當前 time slot 的 active beam 集合，則 beam $b$ 的初始分配功率可表示為
\begin{equation}
\begin{aligned}
P_{s,b}^{(0)}(t)
&=
\frac{N_{s,b}(t)}
{\sum_{b'\in\mathcal{B}_s^{\mathrm{act}}(t)} N_{s,b'}(t)}
\,P_s^{\max},
\end{aligned}
\end{equation}
其中 $b\in\mathcal{B}_s^{\mathrm{act}}(t)$。

為避免將過高功率不切實際地集中於單一 beam，本文進一步考慮每個 beam 的最大功率限制，因此 beam 的最終分配功率表示為
\begin{equation}
P_{s,b}(t)=\min\!\left(P_{s,b}^{(0)}(t),\,P_{s,b}^{\max}\right),
\qquad b\in\mathcal{B}_s^{\mathrm{act}}(t).
\end{equation}

上述配置提供一個 load-aware 的初始功率分配結果。然而，最終 beam power allocation 除需滿足衛星總功率限制與 beam-level 最大功率限制外，仍須進一步滿足對受保護 GSO ground station 的 aggregate EPFD constraint，因此後續仍需執行 power control 或其他資源調整機制。

\subsubsection{SINR and User Satisfaction Model}

在本研究中，user 的接收品質由 serving beam 所提供之期望訊號、同頻干擾以及接收雜訊共同決定。與 full frequency reuse 不同的是，本文採用 OneWeb-like 8-channel frequency assignment。因此，user 所受到的 co-channel interference 並非來自所有 active beams，而僅來自與其 serving beam 使用相同 channel 的 beams。

設 $\mathcal{K}$ 為可用 channel 集合，且 $|\mathcal{K}|=8$。若 user $u$ 由衛星 $s$ 的 beam $b$ 服務，則該 beam 所使用之 channel 記為 $k_{s,b}\in\mathcal{K}$。則 user $u$ 的接收 SINR 可定義為
\begin{equation}
\mathrm{SINR}_u(t)=
\frac{S_u(t)}
{I_u^{\mathrm{intra}}(t)+I_u^{\mathrm{inter}}(t)+N_u(t)}.
\end{equation}

其中 $S_u(t)$ 為期望訊號功率，$I_u^{\mathrm{intra}}(t)$ 與 $I_u^{\mathrm{inter}}(t)$ 分別表示 intra-satellite 同頻干擾與 inter-satellite 同頻干擾，$N_u(t)$ 為接收端雜訊功率。

\paragraph{(1) Desired signal power}
對於由衛星 $s$ 的 beam $b$ 所服務之 user $u$，其期望訊號功率可表示為
\begin{equation}
S_u(t)=
x_{u,s,b}(t)\,
\frac{P_{s,b}(t)\,G^{\mathrm{tx}}_{s,b\to u}(t)\,G_u^{\mathrm{rx}}}
{L_{s,u}(t)}.
\end{equation}

其中 $P_{s,b}(t)$ 為 beam $b$ 的發射功率，$G^{\mathrm{tx}}_{s,b\to u}(t)$ 為 beam 朝 user 方向之發射增益，$G_u^{\mathrm{rx}}$ 為 user 端接收增益，$L_{s,u}(t)$ 為衛星 $s$ 與 user $u$ 之總傳播損失。

本文將總傳播損失表示為
\begin{equation}
L_{s,u}(t)=L_{s,u}^{\mathrm{fs}}(t)\,L_{s,u}^{\mathrm{gas}}(t),
\end{equation}
其中 $L_{s,u}^{\mathrm{fs}}(t)$、$L_{s,u}^{\mathrm{gas}}(t)$ 分別代表 free-space path loss 與 atmospheric gas loss。對於 free-space path loss，其線性形式可寫為
\begin{equation}
L_{s,u}^{\mathrm{fs}}(t)=\left(\frac{4\pi d_{s,u}(t)}{\lambda}\right)^2,
\end{equation}
其中 $d_{s,u}(t)$ 為衛星 $s$ 與 user $u$ 之間的距離，$\lambda$ 為載波波長。

\paragraph{(2) Intra-satellite interference}
由於同一顆衛星上的不同 beams 可能重複使用相同 channel，因此 user 可能受到同一顆 serving satellite 上其他同頻 beams 的干擾。本文將此類干擾定義為 intra-satellite interference，其表達式為
\begin{equation}
\begin{aligned}
I_u^{\mathrm{intra}}(t)
&=
\sum_{\substack{b'\in\mathcal{B}_s\\ b'\neq b}}
\frac{P_{s,b'}(t)\,G^{\mathrm{tx}}_{s,b'\to u}(t)\,G_u^{\mathrm{rx}}}
{L_{s,u}(t)}.
\end{aligned}
\end{equation}

\paragraph{(3) Inter-satellite interference}

除同星同頻干擾外，user 亦可能受到其他可視衛星之同頻 beams 所造成之干擾。本文將此類干擾定義為 inter-satellite interference，其表達式如下。

\begin{equation}
I_u^{\mathrm{inter}}(t)=
\sum_{\substack{s'\in\mathcal{S}(t)\\ s'\neq s}}
\;
\sum_{b'\in\mathcal{B}_{s'}}
\frac{
P_{s',b'}(t)\,G^{\mathrm{tx}}_{s',b'\to u}(t)\,G_u^{\mathrm{rx}}
}{
L_{s',u}(t)
}.
\end{equation}

上式表示，來自其他可視衛星且與 serving beam 使用相同 channel 的 beams，皆可能對 user 造成 co-channel interference。

\paragraph{(4) SINR expression}

綜合上述定義，user 的接收 SINR 可表示如下。

\begin{equation}
\mathrm{SINR}_u(t)=
\frac{S_u(t)}
{I_u^{\mathrm{intra}}(t)+I_u^{\mathrm{inter}}(t)+N_u(t)}.
\end{equation}

上式顯示，在 OneWeb-like 8-channel 架構下，user 的接收品質不僅取決於 serving beam 的功率與幾何位置，亦同時受到 intra-satellite 與 inter-satellite beams 干擾影響。因此，beam power allocation 與 channel-aware user association 將直接影響系統之整體服務表現。

\paragraph{(5) Instantaneous capacity and actual service rate}
在求得 user 的 SINR 後，可依 Shannon formula 計算其 instantaneous capacity：
\begin{equation}
C_u(t)=B_u\log_2\!\bigl(1+\mathrm{SINR}_u(t)\bigr).
\end{equation}

其中 $B_u$ 為 user 所使用之有效頻寬。

由於同一 beam 內可能同時服務多位 users，本文以 TDMA 描述 beam 內資源共享。設 user $u$ 由衛星 $s$ 的 beam $b$ 服務，且該 beam 上共有 $N_{s,b}(t)$ 位 users，則其實際服務速率可表示為
\begin{equation}
R_u(t)=\frac{C_u(t)}{N_{s,b}(t)}.
\end{equation}

\paragraph{(6) User satisfaction}

為評估 user 實際服務品質是否滿足其需求，本文定義 satisfaction 如下。

\begin{equation}
\mathrm{sat}_u(t)=\min\!\left(\frac{R_u(t)}{D_u(t)},\,1\right).
\end{equation}

其中，$D_u(t)$ 為 user $u$ 之需求速率。當 $\mathrm{sat}_u(t)=1$ 時，表示 user 需求已被完全滿足；當 $\mathrm{sat}_u(t)<1$ 時，表示實際服務速率不足以完全支撐其需求。此 satisfaction metric 將作為後續 problem formulation 與 performance evaluation 的主要指標之一。

\subsubsection{Beam-Level EPFD Model for 16-Beam Satellite}

在 NGSO--GSO coexistence 情境下，除需考慮 user link quality 外，亦必須確保 LEO 衛星對受保護 GSO ground station (GS) 所造成之 aggregate EPFD 不超過法規限制。由於本研究中每顆 LEO 衛星配置 16 個 beams，且不同 beams 對 GS 的方向、離軸角與發射功率可能不同，因此單顆衛星對 GS 的干擾不能以單一總功率近似，而應以 beam-level contribution 個別計算後，再於線性域加總。

\paragraph{(1) Beam-level PFD toward the protected GS}
考慮衛星 $s$ 的 beam $b$，其朝受保護 GS 方向所形成之 power flux density (PFD) 可表示為
\begin{equation}
\mathrm{PFD}_{s,b}^{\mathrm{GS}}(t)=
\frac{P_{s,b}(t)\,G^{\mathrm{tx}}_{s,b\to g}(t)}
{4\pi d_{s,g}^{2}(t)}.
\end{equation}

其中 $G^{\mathrm{tx}}_{s,b\to g}(t)$ 為 beam 朝 GS 方向之發射增益，$d_{s,g}(t)$ 為衛星 $s$ 與受保護 GS 之間的距離。

\paragraph{(2) Beam-level EPFD contribution}

EPFD 除考慮到衛星發射至 GS 所形成之 PFD 外，亦需進一步考慮 GS 接收天線對 LEO 方向的離軸接收增益，並相對於其最大接收增益進行正規化。故 beam 對 GS 的 EPFD contribution 可表示如下。

\begin{equation}
\mathrm{EPFD}_{s,b}(t)=
\frac{P_{s,b}(t)\,G^{\mathrm{tx}}_{s,b\to g}(t)}
{4\pi d_{s,g}^{2}(t)}
\cdot
\frac{G_g^{\mathrm{rx}}(\theta_{g\to s}(t))}
{G_{g,\max}^{\mathrm{rx}}}
\cdot
\frac{1}{B_{\mathrm{ref}}}.
\end{equation}

其中，$G_g^{\mathrm{rx}}(\theta_{g\to s}(t))$ 為 GS 天線朝衛星 $s$ 方向之接收增益，$G_{g,\max}^{\mathrm{rx}}$ 為 GS 接收天線最大增益，$B_{\mathrm{ref}}$ 為 EPFD 評估所採用之 reference bandwidth。

\paragraph{(3) Angular dependence of beam gain}
對於衛星 $s$ 的 beam $b$，GS 相對於該 beam boresight 的離軸角記為 $\phi_{s,b\to g}(t)$。因此，beam 在 GS 方向之發射增益可視為離軸角的函數，即
\begin{equation}
G^{\mathrm{tx}}_{s,b\to g}(t)=G^{\mathrm{tx}}_{s,b}\!\bigl(\phi_{s,b\to g}(t)\bigr).
\end{equation}

由於同一顆衛星上的 16 個 beams 各自對應不同覆蓋方向，因此即使它們共享同一顆衛星平台，對 GS 的離軸角仍不相同，進而使各 beam 的 EPFD contribution 顯著不同。這也是本文採用 beam-level EPFD calculation，而非以單一衛星總功率近似的主要原因。

\paragraph{(4) Per-satellite EPFD aggregation over 16 beams}

由於每顆衛星具有 16 個 beams，因此衛星 $s$ 對 GS 的總 EPFD contribution 應為所有 active beams 貢獻之線性總和。設 $\mathcal{B}_s^{\mathrm{act}}(t)$ 為衛星 $s$ 在當前 time slot 的 active beam 集合，則衛星 $s$ 的總 EPFD contribution 可表示如下。

\begin{equation}
\mathrm{EPFD}_{s}^{\mathrm{sat}}(t)=
\sum_{b\in\mathcal{B}_s^{\mathrm{act}}(t)} \mathrm{EPFD}_{s,b}(t).
\end{equation}

由於本研究中每顆衛星固定配置 16 個 beams，因此亦可寫為

\begin{equation}
\mathrm{EPFD}_{s}^{\mathrm{sat}}(t)=
\sum_{b\in\mathcal{B}_s} \mathrm{EPFD}_{s,b}(t),
\end{equation}

其中對於未啟動或未分配功率之 beam，可視為 $P_{s,b}(t)=0$。

\paragraph{(5) Aggregate EPFD from multiple visible satellites}
在當前 time slot 中，所有可視 LEO 衛星對受保護 GS 所形成之 aggregate EPFD 可表示為
\begin{equation}
\mathrm{EPFD}_{\mathrm{agg}}(t)=
\sum_{s\in\mathcal{S}(t)} \mathrm{EPFD}_{s}^{\mathrm{sat}}(t).
\end{equation}

等價地，也可先在單顆衛星內部對 16 個 beams 加總，再對所有可視衛星做加總：
\begin{equation}
\mathrm{EPFD}_{\mathrm{agg}}(t)=
\sum_{s\in\mathcal{S}(t)}\sum_{b\in\mathcal{B}_s} \mathrm{EPFD}_{s,b}(t).
\end{equation}

由於 EPFD 需於線性域累加，因此在進行多 beam、多衛星干擾整合時，不可直接於 dB 域做算術加總，而應先將各 beam contribution 轉為線性值後相加，最後再視需要轉換回 dB 表示。

\paragraph{(6) EPFD constraint}
為確保 NGSO--GSO coexistence 之法規合規性，系統必須滿足下列 aggregate EPFD 限制：
\begin{equation}
\mathrm{EPFD}_{\mathrm{agg}}(t)\le \mathrm{EPFD}_{\mathrm{lim}}.
\end{equation}

其中 $\mathrm{EPFD}_{\mathrm{lim}}$ 為對應場景與 reference bandwidth 下之 EPFD regulatory limit。

當某一 time slot 中 $\mathrm{EPFD}_{\mathrm{agg}}(t)>\mathrm{EPFD}_{\mathrm{lim}}$ 時，表示系統當前 beam power allocation 對受保護 GS 所造成之總干擾超出容許範圍，因此後續必須透過 power control 或其他資源調整機制降低 beam-level contributions，使系統重新滿足法規要求。

\subsection{Problem Formulation}

根據前述系統模型，本文考慮之資源分配問題同時涉及 user association、beam-level power allocation、co-channel interference，以及對受保護 GSO ground station 之 EPFD constraint。研究目標是在滿足 EPFD 法規限制之前提下，盡可能提升整體 user satisfaction。

令系統整體效益函數定義為所有 users satisfaction 之總和，如下所示。

\begin{equation}
\max_{\{x_{u,s,b}(t)\},\{P_{s,b}(t)\}}
\sum_{u\in\mathcal{U}(t)} \mathrm{sat}_u(t).
\end{equation}

則本文所考慮之聯合最佳化問題可表示如下。

\begin{equation}
\begin{aligned}
\max_{\{x_{u,s,b}(t)\},\{P_{s,b}(t)\}} \quad
& \sum_{u\in\mathcal{U}(t)} \mathrm{sat}_u(t)
\end{aligned}
\end{equation}

subject to
\begin{equation}
x_{u,s,b}(t)=0,
\qquad \forall u\notin\mathcal{C}_{s,b}(t),
\end{equation}

\begin{equation}
\sum_{s\in\mathcal{S}(t)}\sum_{b\in\mathcal{B}_s} x_{u,s,b}(t)\le 1,
\qquad \forall u\in\mathcal{U}(t),
\end{equation}

\begin{equation}
\sum_{b\in\mathcal{B}_s} P_{s,b}(t)\le P_s^{\max},
\qquad \forall s\in\mathcal{S}(t),
\end{equation}

\begin{equation}
0\le P_{s,b}(t)\le P_{s,b}^{\max},
\qquad \forall s\in\mathcal{S}(t),\ \forall b\in\mathcal{B}_s,
\end{equation}

\begin{equation}
\mathrm{EPFD}_{\mathrm{agg}}(t)\le \mathrm{EPFD}_{\mathrm{lim}},
\end{equation}

其中，其中，各限制式之物理意義說明如下。

式(27)如果 user u 不在衛星 s 的 beam b 覆蓋範圍內，那這個 association 不允許成立。

式(28)限制每位 user 於每個 time slot 中僅能被一個 satellite-beam 服務。

式(29)限制每顆衛星之總發射功率不得超過其 power budget。

式(30)限制各 beam 之發射功率不得超過 beam-level maximum power。

式(31)確保系統 aggregate EPFD 不超過 regulatory limit。

此目標函數用以在 EPFD 保護限制與衛星資源限制下，最大化可視 LEO 衛星系統中所有 users 的整體 satisfaction。需注意的是，由於 LEO 衛星具有時間變動之幾何關係，因此此最佳化問題需於每一個 time slot 重複求解。

\subsubsection{Discussion on Problem Complexity}

由上述式子可知，本文之最佳化問題屬於一個高複雜度之 mixed-integer non-linear optimization problem。特別是，beam power allocation 會同時影響 user-side SINR 與 GS-side aggregate EPFD，而 user reassociation 又會改變 beam loading，進一步影響 TDMA rate 與 satisfaction。因此，直接求得其全域最優解在實際計算上相當困難。

基於上述原因，本文於下一章提出一套分階段之 heuristic framework，在滿足 EPFD constraint 的前提下，逐步調整 beam powers 與 user reassociation，以提升整體 user satisfaction。
% % Generated from word_IV_v1.docx (Chapter 4 — SAPR-R)
% Include from main.tex via % Generated from word_IV_v1.docx (Chapter 4 — SAPR-R)
% Include from main.tex via % Generated from word_IV_v1.docx (Chapter 4 — SAPR-R)
% Include from main.tex via \input{chapter4}; compile chapter4_standalone.tex for preview.

\input{chapter3_style}

\section{Method}
\label{sec:chapter4}

\subsection{Method Overview}

第三章已將本文問題表述為在 user association、beam power、satellite total power 與 EPFD constraints 下最大化 total user satisfaction 的最佳化問題。然而，該問題同時包含離散的 user--beam association decision 與連續的 beam power configuration，且 association、SINR、TDMA sharing、satisfaction 與 aggregate EPFD 彼此耦合，因此直接求解全域最佳解具有較高計算複雜度。

為降低 online decision complexity，本文提出一套 Safe-Beam Assisted Power-Release Relay (SAPR-R) 方法。SAPR-R 利用 LEO constellation、protected GS location、GSO direction 與 fixed beam layout 的可預測性，先針對整個觀察期間的所有 time slots 建立 beam-level relation records。對於每個 time slot，系統預先記錄 critical satellite、north/south helper satellites、critical satellite 上需關閉的 beams、仍可運作的 safe beams，以及 critical beams 與 helper beams 之間的 beam-pair mappings。此階段僅依賴 satellite geometry、beam footprint 與 EPFD calculation，不需要 user distribution，因此可於 online user reassociation 前完成。

在 online operation 階段，系統讀取當前 time slot 的 beam relation record，並根據當下 users 的位置與 association 執行 user-aware relay decision。首先，SAPR-R 透過 Safe-Beam Assisted Helper Relief (SAHR) 使 critical safe beams 嘗試接手 helper release beams 上的部分 users。若 reassociation 不降低被轉移 user 與 safe beam 原本 users 的 satisfaction，則接受此次轉移，並將 helper release beam 所釋放出的功率回收到 helper satellite 的 available power pool。接著，helper satellite 依據 relay demand 將 power pool 分配給 helper relay beams，以提升其後續接手 closed-beam users 的 capacity。最後，SAPR-R 透過 Beam-Pair Local Relay (BPLR)，在每組 closed beam/helper relay beam pair 內，依據 relay link quality 對 critical closed-beam users 進行排序與 sequential allocation。

綜合而言，本文方法包含下列步驟：
\begin{itemize}[leftmargin=*,nosep]
  \item \textbf{Time-slot beam relation record:} 預先記錄每個 time slot 的 critical/helper satellites、closed/safe beams 與 beam-pair mappings。
  \item \textbf{Safe-Beam Assisted Helper Relief (SAHR):} 利用 critical safe beams 接手部分 helper users，在不降低 satisfaction 的條件下釋放 helper satellite power。
  \item \textbf{Helper relay beam boosting:} 依據 relay demand 將 available power pool 分配給 helper relay beams。
  \item \textbf{Beam-Pair Local Relay (BPLR):} 在每組 beam pair 內，依 relay link quality 分配 relay capacity 給 closed-beam users。
\end{itemize}

\subsection{Time-Slot Beam Relation Record}

\begin{figure*}[t]
  \centering
  \includegraphics[width=\textwidth]{ch4_beam_relation_record}
  \caption{三個 time slots 中，沿軌道方向之三顆 LEO 衛星（$s^{\mathrm{N}}$、$s^{\mathrm{c}}$、$s^{\mathrm{S}}$）飛過受保護 GSO 地面站之示意。critical 衛星區域 $S_{\mathrm{C}}$ 內，黑色斜線為 $\mathcal{B}_c^{\mathrm{off}}(t)$（closed beams），綠色斜線為 $\mathcal{B}_c^{\mathrm{safe}}(t)$（safe beams）。(a)--(c) 為由早至晚之 time slot。}
  \label{fig:ch4_beam_relation_record}
\end{figure*}

對於每一個 time slot $t$，本文建立一個 beam relation record $\mathcal{R}_B(t)$。圖~\ref{fig:ch4_beam_relation_record} 示意本文所考慮之 time-slot beam 狀態，包括：
\begin{equation*}
\begin{split}
\mathcal{R}_B(t)=\bigl(
& s^{\mathrm{c}}(t),\,
s^{\mathrm{N}}(t),\,
s^{\mathrm{S}}(t), \\
& \mathcal{B}_c^{\mathrm{off}}(t),\,
\mathcal{B}_c^{\mathrm{safe}}(t), \\
& \mathcal{M}_{\mathrm{safe}\to\mathrm{rel}}(t),\,
\mathcal{M}_{\mathrm{off}\to\mathrm{relay}}(t)
\bigr).
\end{split}
\end{equation*}
該 record 儲存後續 SAHR 與 BPLR 所需的 beam-level information，其中 $s^{\mathrm{c}}(t)$、$s^{\mathrm{N}}(t)$ 與 $s^{\mathrm{S}}(t)$ 分別為 critical satellite、north helper satellite 與 south helper satellite。$\mathcal{B}_c^{\mathrm{off}}(t)$ 為 critical satellite 上因 EPFD constraint 需被關閉的 beams，$\mathcal{B}_c^{\mathrm{safe}}(t)$ 為未被關閉且維持 baseline power 的 safe beams。$\mathcal{M}_{\mathrm{safe}\to\mathrm{rel}}(t)$ 與 $\mathcal{M}_{\mathrm{off}\to\mathrm{relay}}(t)$ 分別表示 safe-beam-to-helper-release-beam mapping 與 closed-beam-to-helper-relay-beam mapping。

\mdsubsubsection{Critical and Helper Satellite Identification}
在每個 time slot $t$，系統根據 LEO satellite 與 protected GS 之間的幾何關係，找出對 GS 造成主要 EPFD 風險的 critical satellite：
\begin{equation*}
s^{\mathrm{c}}(t).
\end{equation*}
根據模擬觀察，EPFD violation 主要由 critical satellite 上的 high-risk beams 主導，因此本文僅針對 $s^{\mathrm{c}}(t)$ 上的 beams 進行 shutdown decision，不關閉 helper satellites 或其他 visible satellites 的 beams。接著，選取 critical satellite 同軌道方向上的 north neighbor 與 south neighbor 作為 helper satellites：
\begin{equation*}
s^{\mathrm{N}}(t),\quad s^{\mathrm{S}}(t).
\end{equation*}

\mdsubsubsection{EPFD-Driven Beam Status Determination}
對每個 critical beam $b\in\mathcal{B}_{s^{\mathrm{c}}}$，定義其 EPFD risk score 為
\begin{equation*}
q_{c,b}(t)=\mathrm{EPFD}_{s^{\mathrm{c}},b}(t).
\end{equation*}
critical beams 依據 $q_{c,b}(t)$ 由大到小排序，並採用 greedy shutdown rule：每次關閉目前 EPFD risk score 最高的 beam，直到 aggregate EPFD 滿足
\begin{equation*}
\mathrm{EPFD}_{\mathrm{agg}}(t)\le \mathrm{EPFD}_{\mathrm{lim}}.
\end{equation*}
被關閉的 beams 組成
\begin{equation*}
\mathcal{B}_c^{\mathrm{off}}(t),
\end{equation*}
其餘 beams 則形成 safe beam set：
\begin{equation*}
\mathcal{B}_c^{\mathrm{safe}}(t)=\mathcal{B}_{s^{\mathrm{c}}}\setminus \mathcal{B}_c^{\mathrm{off}}(t).
\end{equation*}
此階段僅根據 beam-level EPFD contribution 與 regulatory limit 決定 beam status，不涉及 user distribution，因此可於 online user reassociation 前預先完成。

\mdsubsubsection{Beam-Pair Mapping Construction}
在得到 $\mathcal{B}_c^{\mathrm{off}}(t)$ 與 $\mathcal{B}_c^{\mathrm{safe}}(t)$ 後，本文根據 beam footprint overlap 建立 beam-pair mappings。對於 SAHR，定義
\begin{equation*}
\begin{split}
\mathcal{M}_{\mathrm{safe}\to\mathrm{rel}}(t)=
\bigl\{
&(b_c^{\mathrm{safe}},b_h^{\mathrm{rel}}) : \\
&\mathcal{F}_{s^{\mathrm{c}},b_c^{\mathrm{safe}}}(t)\cap
\mathcal{F}_{s_h,b_h^{\mathrm{rel}}}(t)\neq\emptyset, \\
& s_h\in\{s^{\mathrm{N}}(t),s^{\mathrm{S}}(t)\}
\bigr\},
\end{split}
\end{equation*}
其中 $s_h\in\{s^{\mathrm{N}}(t),s^{\mathrm{S}}(t)\}$。此 mapping 用於找出 critical safe beam 可接手哪些 helper release beam 上的 users。

對於 BPLR，定義
\begin{equation*}
\begin{split}
\mathcal{M}_{\mathrm{off}\to\mathrm{relay}}(t)=
\bigl\{
&(b_c^{\mathrm{off}},b_h^{\mathrm{relay}}) : \\
&\mathcal{F}_{s^{\mathrm{c}},b_c^{\mathrm{off}}}(t)\cap
\mathcal{F}_{s_h,b_h^{\mathrm{relay}}}(t)\neq\emptyset
\bigr\}.
\end{split}
\end{equation*}
此 mapping 用於找出 critical closed-beam users 可由哪些 helper relay beams 接手。

圖~\ref{fig:ch4_beam_pair_mapping} 呈現儲存於 $\mathcal{R}_B(t)$ 中的兩組 beam-pair mappings：$\mathcal{M}_{\mathrm{safe}\to\mathrm{rel}}(t)$（SAHR）與 $\mathcal{M}_{\mathrm{off}\to\mathrm{relay}}(t)$（BPLR）。

\begin{figure*}[t]
  \centering
  \includegraphics[width=\textwidth]{ch4_beam_pair_mapping}
  \caption{Time slot $t$ 之 beam-pair mapping 對照表。(a) $\mathcal{M}_{\mathrm{safe}\to\mathrm{rel}}(t)$：$s^{\mathrm{c}}$ 上 safe beams 對應至 $s^{\mathrm{N}}$/$s^{\mathrm{S}}$ 之 helper release beams（SAHR）。(b) $\mathcal{M}_{\mathrm{off}\to\mathrm{relay}}(t)$：closed beams 對應至 helper relay beams（BPLR）。}
  \label{fig:ch4_beam_pair_mapping}
\end{figure*}

本文考慮的 Walker-star constellation 與 fixed beam layout 使 critical satellite 與 north/south helper satellites 呈現近似 one-to-one beam overlap。具體而言，critical satellite 的 north-half beams 主要與 north helper satellite 的 south-half beams 重疊，而 critical satellite 的 south-half beams 主要與 south helper satellite 的 north-half beams 重疊。因此，後續 relay 可在 beam pair 內局部執行，而不需進行 global matching。

此外，helper beams 分為 helper release beams 與 helper relay beams。前者用於釋放 helper satellite power，後者用於接手 critical closed-beam users。本文假設兩者於同一 time slot 中不重疊：
\begin{equation*}
\mathcal{B}_h^{\mathrm{rel}}(t)\cap \mathcal{B}_h^{\mathrm{relay}}(t)=\emptyset.
\end{equation*}

\subsection{Safe-Beam Assisted Helper Relief}

讀取 $\mathcal{R}_B(t)$ 後，online 程序首先執行 SAHR。圖~\ref{fig:ch4_Safe-beam_assisted_power_release} 示意 SAHR 流程。此階段的目的不是直接接手 critical closed-beam users，而是利用 critical safe beams 接手部分 helper users，使 helper satellite 釋放功率，供後續 relay boosting 使用。

對於
\begin{equation*}
(b_c^{\mathrm{safe}},b_h^{\mathrm{rel}})\in \mathcal{M}_{\mathrm{safe}\to\mathrm{rel}}(t),
\end{equation*}
系統找出目前由 $b_h^{\mathrm{rel}}$ 服務，且同時位於 $b_c^{\mathrm{safe}}$ footprint 內的 helper users：
\begin{equation*}
\begin{split}
&\mathcal{V}_{\mathrm{rel}}(b_c^{\mathrm{safe}},b_h^{\mathrm{rel}},t)= \\
&\quad
\left\{
v :
v\in\mathcal{U}_{s_h,b_h^{\mathrm{rel}}}(t),\;
v\in\mathcal{F}_{s^{\mathrm{c}},b_c^{\mathrm{safe}}}(t)
\right\}.
\end{split}
\end{equation*}
若將 candidate helper user $v$ reassociate 至 $b_c^{\mathrm{safe}}$，helper release beam 所需功率的下降量為
\begin{equation*}
\Delta P_v^{\mathrm{rel}}(t)=
P_{b_h^{\mathrm{rel}}}^{\mathrm{req,old}}(t)-
P_{b_h^{\mathrm{rel}}}^{\mathrm{req,new}}(v,t).
\end{equation*}
SAHR score 定義為
\begin{equation*}
S_{\mathrm{relief}}(v,t)=\Delta P_v^{\mathrm{rel}}(t).
\end{equation*}
candidate helper users 依據 $S_{\mathrm{relief}}(v,t)$ 由大到小排序。僅當被轉移 user 的 satisfaction 不下降，且 safe beam 原本 users 亦不下降時，才接受 reassociation：
\begin{equation*}
\mathrm{sat}_v^{\mathrm{safe}}(t)\ge \mathrm{sat}_v^{\mathrm{old}}(t),
\end{equation*}
\begin{equation*}
\mathrm{sat}_w^{\mathrm{new}}(t)\ge \mathrm{sat}_w^{\mathrm{old}}(t),
\qquad \forall w\in\mathcal{U}_{b_c^{\mathrm{safe}}}^{\mathrm{ori}}(t).
\end{equation*}
若條件成立，user $v$ 被 reassociate 至 critical safe beam，且 $\Delta P_v^{\mathrm{rel}}(t)$ 回到 helper satellite power pool。對於 helper satellite $s_h$，SAHR 釋放的總功率為
\begin{equation*}
P_{s_h}^{\mathrm{rel}}(t)=
\sum_{v\in\mathcal{A}_{\mathrm{rel}}(s_h,t)} \Delta P_v^{\mathrm{rel}}(t),
\end{equation*}
其中 $\mathcal{A}_{\mathrm{rel}}(s_h,t)$ 為成功被 safe beams 接手的 helper users 集合。

baseline operation 下原本未使用的功率為
\begin{equation*}
P_{s_h}^{\mathrm{unused}}(t)=
P_{s_h}^{\max}-\sum_{b\in\mathcal{B}_{s_h}} P^{\mathrm{base}},
\end{equation*}
因此，後續可用於 relay boosting 的 power pool 為
\begin{equation*}
P_{s_h}^{\mathrm{pool}}(t)=
P_{s_h}^{\mathrm{unused}}(t)+P_{s_h}^{\mathrm{rel}}(t).
\end{equation*}
除 helper release beams 與 helper relay beams 外，其餘 helper beams 皆維持 baseline power $P^{\mathrm{base}}$。

\begin{figure*}[t]
  \centering
  \includegraphics[width=\textwidth]{ch4_Safe-beam_assisted_power_release}
  \caption{Safe-Beam Assisted Helper Relief (SAHR) 示意。critical safe beams 接手 helper release beams 上符合條件的 users，所釋放功率回收到 helper satellite 的 power pool。}
  \label{fig:ch4_Safe-beam_assisted_power_release}
\end{figure*}

\subsection{Helper Relay Beam Boosting}

完成 SAHR 後，helper satellite $s_h$ 使用 $P_{s_h}^{\mathrm{pool}}(t)$ 提升 helper relay beams 的 transmit power。helper relay beams 由 $\mathcal{M}_{\mathrm{off}\to\mathrm{relay}}(t)$ 決定：
\begin{equation*}
\begin{split}
\mathcal{B}_{s_h}^{\mathrm{relay}}(t)=
\bigl\{
&b_h^{\mathrm{relay}} :
\exists\, b_c^{\mathrm{off}}\in\mathcal{B}_c^{\mathrm{off}}(t): \\
&\qquad (b_c^{\mathrm{off}},b_h^{\mathrm{relay}})\in \mathcal{M}_{\mathrm{off}\to\mathrm{relay}}(t)
\bigr\}.
\end{split}
\end{equation*}
將 relay beam 從 baseline power 提升至 beam upper bound 所需的額外功率為
\begin{equation*}
\begin{split}
\Delta P_{b_h^{\mathrm{relay}}}^{\max}(t)=
P_{s_h,b_h^{\mathrm{relay}}}^{\max}-P^{\mathrm{base}}.
\end{split}
\end{equation*}
對於每個 helper relay beam，其 relay demand 定義為
\begin{equation*}
D_{b_h^{\mathrm{relay}}}^{\mathrm{relay}}(t)=
\left|\mathcal{U}_{\mathrm{cand}}(b_h^{\mathrm{relay}},t)\right|,
\end{equation*}
其中 $\mathcal{U}_{\mathrm{cand}}(b_h^{\mathrm{relay}},t)$ 表示根據 $\mathcal{M}_{\mathrm{off}\to\mathrm{relay}}(t)$ 得到的 candidate critical users，也就是位於對應 critical closed beam 內，且可被 $b_h^{\mathrm{relay}}$ 覆蓋的 users。

若 $P_{s_h}^{\mathrm{pool}}(t)$ 足以支援所有 relay beams 提升至 upper bound，則所有 relay beams 皆可被 boost。若 power pool 不足，本文採用 relay-demand-priority boosting：helper relay beams 依據 $D_{b_h^{\mathrm{relay}}}^{\mathrm{relay}}(t)$ 由大到小排序，relay demand 較高者優先取得額外功率。

在分配過程中，若剩餘 power pool 足以使當前 relay beam 提升至 beam upper bound $P_{s_h,b_h^{\mathrm{relay}}}^{\max}$，則將其提升至上限。若剩餘 power pool 不足以完整提升下一個 relay beam，則仍可將剩餘功率分配給該 beam，使其 transmit power 高於 $P^{\mathrm{base}}$。如此可避免剩餘 power pool 被浪費，並提升該 relay beam 的可用 relay capacity。

因此，所有 helper relay beams 仍可進入後續 BPLR 階段，差別在於各 relay beams 依據最終 transmit power 具有不同的 available relay capacity。

\subsection{Beam-Pair Local Relay}

完成 relay-beam boosting 後，系統於每一組 closed-beam-to-helper-relay-beam pair 內執行 BPLR。圖~\ref{fig:ch4_Beam-pair_local_relay} 示意 BPLR 流程。對於
\begin{equation*}
(b_c^{\mathrm{off}},b_h^{\mathrm{relay}})\in \mathcal{M}_{\mathrm{off}\to\mathrm{relay}}(t),
\end{equation*}
系統找出原本由 $b_c^{\mathrm{off}}$ 服務，且同時位於 $b_h^{\mathrm{relay}}$ footprint 內的 critical users：
\begin{equation*}
\begin{split}
&\mathcal{U}_{\mathrm{pair}}(b_c^{\mathrm{off}},b_h^{\mathrm{relay}},t)= \\
&\quad
\left\{
u :
u\in\mathcal{U}_{s^{\mathrm{c}},b_c^{\mathrm{off}}}(t),\;
u\in\mathcal{F}_{s_h,b_h^{\mathrm{relay}}}(t)
\right\}.
\end{split}
\end{equation*}
對於 candidate critical user $u$，其由 helper relay beam 服務時的預期 achievable rate 記為
\begin{equation*}
R_{u,b_h^{\mathrm{relay}}}(t).
\end{equation*}
本文定義 BPLR score 為
\begin{equation*}
S_{\mathrm{relay}}(u,t)=R_{u,b_h^{\mathrm{relay}}}(t).
\end{equation*}
candidate users 依據 $S_{\mathrm{relay}}(u,t)$ 由大到小排序。排序後，helper relay beam 依據其 final transmit power 計算可用於 critical users 的 remaining relay capacity。此 capacity 已扣除保護原本 helper users 所需的服務資源。

接著，系統依序將 remaining relay capacity 分配給 sorted critical users。若剩餘容量足以滿足 user demand $D_u(t)$，則將該 user reassociate 至 helper relay beam，並更新其 satisfaction。若最後剩餘容量不足以完全滿足下一個 user，但仍大於零，則將剩餘容量分配給該 user，使其獲得 partial satisfaction。當 remaining relay capacity 用盡後，該 beam pair 的 relay process 即停止。

所有 relay decisions 仍需滿足 single association constraint。也就是說，被 relay 的 user 會從 critical closed beam reassociate 至 helper relay beam，且同一 time slot 中仍只由一個 beam 服務。

\begin{figure*}[t]
  \centering
  \includegraphics[width=\textwidth]{ch4_Beam-pair_local_relay}
  \caption{Beam-Pair Local Relay (BPLR) 示意。依 $\mathcal{M}_{\mathrm{off}\to\mathrm{relay}}(t)$ 將 closed beams 上的 candidate critical users reassociate 至對應的 helper relay beams。}
  \label{fig:ch4_Beam-pair_local_relay}
\end{figure*}

\subsection{Overall Algorithm}

本文之 online procedure 採用單一 time-slot-based algorithm。在每個 time slot 中，SAPR-R 先執行 SAHR 與 relay-beam boosting，接著根據更新後的 helper relay beam power 執行 BPLR。完整流程如 Algorithm~\ref{alg:ch4_sapr} 所示。

\begin{algorithm}[H]
\caption{SAPR-R Online Operation}
\label{alg:ch4_sapr}
\scriptsize
\begin{algorithmic}[1]
\Require $\{\mathcal{R}_B(t)\}_{t=1}^{T}$, initial $\mathbf{x}(t)$, $\mathcal{U}(t)$, $D_u(t)$
\Require $P^{\mathrm{base}}$, $P_{s,b}^{\max}$, $P_s^{\max}$, coverage, capacity/satisfaction model
\Ensure $\mathbf{x}^{\mathrm{final}}(t)$, $P_{s,b}^{\mathrm{final}}(t)$, $\mathrm{sat}_u^{\mathrm{final}}(t)$

\For{$t=1$ \textbf{to} $T$}
    \State Load $\mathcal{R}_B(t)$; initialize beam powers per record and baseline operation
    \State \textbf{Stage~1: SAHR and relay-beam boosting}
    \ForAll{$(b_c^{\mathrm{safe}},b_h^{\mathrm{rel}})\in \mathcal{M}_{\mathrm{safe}\to\mathrm{rel}}(t)$}
        \State Build candidate helper users; compute $\Delta P_v^{\mathrm{rel}}(t)$
        \State Sort candidates by $\Delta P_v^{\mathrm{rel}}(t)$ descending
        \ForAll{$v$ in sorted order}
            \State Trial reassociation $v$: $b_h^{\mathrm{rel}}\to b_c^{\mathrm{safe}}$
            \If{$v$ not degraded and users on $b_c^{\mathrm{safe}}$ not degraded}
                \State Accept; update association and $P_{s_h}^{\mathrm{pool}}(t)$
            \EndIf
        \EndFor
    \EndFor
    \ForAll{$s_h\in\{s^{\mathrm{N}}(t),s^{\mathrm{S}}(t)\}$}
        \State Extract $\mathcal{B}_{s_h}^{\mathrm{relay}}(t)$; compute relay demand; sort by demand
        \ForAll{$b_h^{\mathrm{relay}}\in \mathcal{B}_{s_h}^{\mathrm{relay}}(t)$ in sorted order}
            \State Allocate $P_{s_h}^{\mathrm{pool}}(t)$ (full boost, partial boost, or baseline)
        \EndFor
        \State Check $\sum_b P_{s,b}(t)\le P_{s_h}^{\max}$
    \EndFor
    \State \textbf{Stage~2: Beam-pair local relay}
    \ForAll{$(b_c^{\mathrm{off}},b_h^{\mathrm{relay}})\in \mathcal{M}_{\mathrm{off}\to\mathrm{relay}}(t)$}
        \State Build candidate critical users; sort by relay link quality
        \State Compute remaining relay capacity of $b_h^{\mathrm{relay}}$ (protect helper users)
        \ForAll{$u$ in sorted order}
            \If{capacity $\ge D_u(t)$}
                \State Reassociate $u$ to $b_h^{\mathrm{relay}}$; update $\mathrm{sat}_u(t)$ and capacity
            \ElsIf{remaining capacity $>0$}
                \State Partial allocation to $u$; set remaining capacity to $0$; \textbf{break}
            \Else
                \State \textbf{break}
            \EndIf
        \EndFor
    \EndFor
    \State Validate single association, beam power, satellite total power, and EPFD constraints
    \State Store $\mathbf{x}^{\mathrm{final}}(t)$, $P_{s,b}^{\mathrm{final}}(t)$, $\mathrm{sat}_u^{\mathrm{final}}(t)$
\EndFor
\end{algorithmic}
\end{algorithm}

\subsection{Complexity Discussion}

本文方法避免直接求解 mixed-integer nonlinear optimization problem。由於每個 time slot 的 beam relation record 已預先建立，online procedure 不需重新搜尋 critical satellite、helper satellites 或 beam-pair mappings，而是直接在對應 beam pairs 內執行局部 reassociation 與 relay decision。

令 $K_t^{\mathrm{rel}}$ 表示 time slot $t$ 中 SAHR 階段的 candidate helper users 數量，$K_t^{\mathrm{relay}}$ 表示 BPLR 階段的 candidate critical users 數量。在 Algorithm~\ref{alg:ch4_sapr} 中，candidate 建立與 released power 計算皆與 $K_t^{\mathrm{rel}}$ 呈線性關係，而主要成本來自根據 $\Delta P_v^{\mathrm{rel}}(t)$ 對 candidate helper users 排序。因此，SAHR 的時間複雜度為
\begin{equation*}
\mathcal{O}\!\bigl(K_t^{\mathrm{rel}}\log K_t^{\mathrm{rel}}\bigr).
\end{equation*}
在 BPLR 階段，candidate critical users 依據 relay link quality 排序，後續 capacity allocation 則為 sequential allocation。因此，BPLR 的時間複雜度為
\begin{equation*}
\mathcal{O}\!\bigl(K_t^{\mathrm{relay}}\log K_t^{\mathrm{relay}}\bigr).
\end{equation*}
綜合所有 time slots，本文方法的 online complexity 可表示為
\begin{equation*}
\mathcal{O}\!\left(
\sum_{t=1}^{T}
K_t^{\mathrm{rel}}\log K_t^{\mathrm{rel}}
+
K_t^{\mathrm{relay}}\log K_t^{\mathrm{relay}}
\right).
\end{equation*}
此結果顯示，本文方法將 online decision 限制於 precomputed beam pairs 內的 local sorting 與 sequential allocation，而非每個 time slot 重新執行 global user association 與 continuous power optimization。

\subsection{Summary}

本章提出 Safe-Beam Assisted Power-Release Relay (SAPR-R) 方法，以處理 critical satellite 接近 protected GS 時 EPFD constraint 與 user satisfaction 之間的 tradeoff。本文先利用 predictable satellite geometry 與 beam footprint overlap 建立 time-slot beam relation record，再於 online 階段執行 SAHR、relay beam boosting 與 BPLR。

透過此分階段設計，本文將原本高度耦合的 user association 與 beam power allocation 問題拆解為可局部處理的 decision steps。critical safe beams 用於釋放 helper satellite power，helper relay beams 則根據 relay demand 取得額外 transmit power，並於 beam pair 內接手 critical closed-beam users。此方法可在滿足 EPFD constraint 的前提下，提高 critical period 中的 total user satisfaction。
; compile chapter4_standalone.tex for preview.

% Reference style (like IEEE paper): **Label:** regular body text on the same line.
% - Label before colon: bold (bfseries)
% - Following text: normal weight (mdseries / normalfont)
\makeatletter
\@ifundefined{md@runinlabel}{%
  \newcommand{\md@runinlabel}[1]{{\bfseries #1}\normalfont}%
  \newcommand{\mdsubsubsection}[2][]{%
    \refstepcounter{subsubsection}%
    \par\vspace{0.35ex}%
    \noindent\md@runinlabel{\arabic{subsubsection})~#2:}\space
    \def\@mdss@label{#1}%
    \ifx\@mdss@label\@empty\else\label{#1}\fi
  }%
  \newcommand{\mdparagraph}[2][]{%
    \par\vspace{0.25ex}%
    \noindent\md@runinlabel{#2:}\space
    \def\@mdpar@label{#1}%
    \ifx\@mdpar@label\@empty\else\label{#1}\fi
  }%
}{}
\makeatother


\section{Method}
\label{sec:chapter4}

\subsection{Method Overview}

第三章已將本文問題表述為在 user association、beam power、satellite total power 與 EPFD constraints 下最大化 total user satisfaction 的最佳化問題。然而，該問題同時包含離散的 user--beam association decision 與連續的 beam power configuration，且 association、SINR、TDMA sharing、satisfaction 與 aggregate EPFD 彼此耦合，因此直接求解全域最佳解具有較高計算複雜度。

為降低 online decision complexity，本文提出一套 Safe-Beam Assisted Power-Release Relay (SAPR-R) 方法。SAPR-R 利用 LEO constellation、protected GS location、GSO direction 與 fixed beam layout 的可預測性，先針對整個觀察期間的所有 time slots 建立 beam-level relation records。對於每個 time slot，系統預先記錄 critical satellite、north/south helper satellites、critical satellite 上需關閉的 beams、仍可運作的 safe beams，以及 critical beams 與 helper beams 之間的 beam-pair mappings。此階段僅依賴 satellite geometry、beam footprint 與 EPFD calculation，不需要 user distribution，因此可於 online user reassociation 前完成。

在 online operation 階段，系統讀取當前 time slot 的 beam relation record，並根據當下 users 的位置與 association 執行 user-aware relay decision。首先，SAPR-R 透過 Safe-Beam Assisted Helper Relief (SAHR) 使 critical safe beams 嘗試接手 helper release beams 上的部分 users。若 reassociation 不降低被轉移 user 與 safe beam 原本 users 的 satisfaction，則接受此次轉移，並將 helper release beam 所釋放出的功率回收到 helper satellite 的 available power pool。接著，helper satellite 依據 relay demand 將 power pool 分配給 helper relay beams，以提升其後續接手 closed-beam users 的 capacity。最後，SAPR-R 透過 Beam-Pair Local Relay (BPLR)，在每組 closed beam/helper relay beam pair 內，依據 relay link quality 對 critical closed-beam users 進行排序與 sequential allocation。

綜合而言，本文方法包含下列步驟：
\begin{itemize}[leftmargin=*,nosep]
  \item \textbf{Time-slot beam relation record:} 預先記錄每個 time slot 的 critical/helper satellites、closed/safe beams 與 beam-pair mappings。
  \item \textbf{Safe-Beam Assisted Helper Relief (SAHR):} 利用 critical safe beams 接手部分 helper users，在不降低 satisfaction 的條件下釋放 helper satellite power。
  \item \textbf{Helper relay beam boosting:} 依據 relay demand 將 available power pool 分配給 helper relay beams。
  \item \textbf{Beam-Pair Local Relay (BPLR):} 在每組 beam pair 內，依 relay link quality 分配 relay capacity 給 closed-beam users。
\end{itemize}

\subsection{Time-Slot Beam Relation Record}

\begin{figure*}[t]
  \centering
  \includegraphics[width=\textwidth]{ch4_beam_relation_record}
  \caption{三個 time slots 中，沿軌道方向之三顆 LEO 衛星（$s^{\mathrm{N}}$、$s^{\mathrm{c}}$、$s^{\mathrm{S}}$）飛過受保護 GSO 地面站之示意。critical 衛星區域 $S_{\mathrm{C}}$ 內，黑色斜線為 $\mathcal{B}_c^{\mathrm{off}}(t)$（closed beams），綠色斜線為 $\mathcal{B}_c^{\mathrm{safe}}(t)$（safe beams）。(a)--(c) 為由早至晚之 time slot。}
  \label{fig:ch4_beam_relation_record}
\end{figure*}

對於每一個 time slot $t$，本文建立一個 beam relation record $\mathcal{R}_B(t)$。圖~\ref{fig:ch4_beam_relation_record} 示意本文所考慮之 time-slot beam 狀態，包括：
\begin{equation*}
\begin{split}
\mathcal{R}_B(t)=\bigl(
& s^{\mathrm{c}}(t),\,
s^{\mathrm{N}}(t),\,
s^{\mathrm{S}}(t), \\
& \mathcal{B}_c^{\mathrm{off}}(t),\,
\mathcal{B}_c^{\mathrm{safe}}(t), \\
& \mathcal{M}_{\mathrm{safe}\to\mathrm{rel}}(t),\,
\mathcal{M}_{\mathrm{off}\to\mathrm{relay}}(t)
\bigr).
\end{split}
\end{equation*}
該 record 儲存後續 SAHR 與 BPLR 所需的 beam-level information，其中 $s^{\mathrm{c}}(t)$、$s^{\mathrm{N}}(t)$ 與 $s^{\mathrm{S}}(t)$ 分別為 critical satellite、north helper satellite 與 south helper satellite。$\mathcal{B}_c^{\mathrm{off}}(t)$ 為 critical satellite 上因 EPFD constraint 需被關閉的 beams，$\mathcal{B}_c^{\mathrm{safe}}(t)$ 為未被關閉且維持 baseline power 的 safe beams。$\mathcal{M}_{\mathrm{safe}\to\mathrm{rel}}(t)$ 與 $\mathcal{M}_{\mathrm{off}\to\mathrm{relay}}(t)$ 分別表示 safe-beam-to-helper-release-beam mapping 與 closed-beam-to-helper-relay-beam mapping。

\mdsubsubsection{Critical and Helper Satellite Identification}
在每個 time slot $t$，系統根據 LEO satellite 與 protected GS 之間的幾何關係，找出對 GS 造成主要 EPFD 風險的 critical satellite：
\begin{equation*}
s^{\mathrm{c}}(t).
\end{equation*}
根據模擬觀察，EPFD violation 主要由 critical satellite 上的 high-risk beams 主導，因此本文僅針對 $s^{\mathrm{c}}(t)$ 上的 beams 進行 shutdown decision，不關閉 helper satellites 或其他 visible satellites 的 beams。接著，選取 critical satellite 同軌道方向上的 north neighbor 與 south neighbor 作為 helper satellites：
\begin{equation*}
s^{\mathrm{N}}(t),\quad s^{\mathrm{S}}(t).
\end{equation*}

\mdsubsubsection{EPFD-Driven Beam Status Determination}
對每個 critical beam $b\in\mathcal{B}_{s^{\mathrm{c}}}$，定義其 EPFD risk score 為
\begin{equation*}
q_{c,b}(t)=\mathrm{EPFD}_{s^{\mathrm{c}},b}(t).
\end{equation*}
critical beams 依據 $q_{c,b}(t)$ 由大到小排序，並採用 greedy shutdown rule：每次關閉目前 EPFD risk score 最高的 beam，直到 aggregate EPFD 滿足
\begin{equation*}
\mathrm{EPFD}_{\mathrm{agg}}(t)\le \mathrm{EPFD}_{\mathrm{lim}}.
\end{equation*}
被關閉的 beams 組成
\begin{equation*}
\mathcal{B}_c^{\mathrm{off}}(t),
\end{equation*}
其餘 beams 則形成 safe beam set：
\begin{equation*}
\mathcal{B}_c^{\mathrm{safe}}(t)=\mathcal{B}_{s^{\mathrm{c}}}\setminus \mathcal{B}_c^{\mathrm{off}}(t).
\end{equation*}
此階段僅根據 beam-level EPFD contribution 與 regulatory limit 決定 beam status，不涉及 user distribution，因此可於 online user reassociation 前預先完成。

\mdsubsubsection{Beam-Pair Mapping Construction}
在得到 $\mathcal{B}_c^{\mathrm{off}}(t)$ 與 $\mathcal{B}_c^{\mathrm{safe}}(t)$ 後，本文根據 beam footprint overlap 建立 beam-pair mappings。對於 SAHR，定義
\begin{equation*}
\begin{split}
\mathcal{M}_{\mathrm{safe}\to\mathrm{rel}}(t)=
\bigl\{
&(b_c^{\mathrm{safe}},b_h^{\mathrm{rel}}) : \\
&\mathcal{F}_{s^{\mathrm{c}},b_c^{\mathrm{safe}}}(t)\cap
\mathcal{F}_{s_h,b_h^{\mathrm{rel}}}(t)\neq\emptyset, \\
& s_h\in\{s^{\mathrm{N}}(t),s^{\mathrm{S}}(t)\}
\bigr\},
\end{split}
\end{equation*}
其中 $s_h\in\{s^{\mathrm{N}}(t),s^{\mathrm{S}}(t)\}$。此 mapping 用於找出 critical safe beam 可接手哪些 helper release beam 上的 users。

對於 BPLR，定義
\begin{equation*}
\begin{split}
\mathcal{M}_{\mathrm{off}\to\mathrm{relay}}(t)=
\bigl\{
&(b_c^{\mathrm{off}},b_h^{\mathrm{relay}}) : \\
&\mathcal{F}_{s^{\mathrm{c}},b_c^{\mathrm{off}}}(t)\cap
\mathcal{F}_{s_h,b_h^{\mathrm{relay}}}(t)\neq\emptyset
\bigr\}.
\end{split}
\end{equation*}
此 mapping 用於找出 critical closed-beam users 可由哪些 helper relay beams 接手。

圖~\ref{fig:ch4_beam_pair_mapping} 呈現儲存於 $\mathcal{R}_B(t)$ 中的兩組 beam-pair mappings：$\mathcal{M}_{\mathrm{safe}\to\mathrm{rel}}(t)$（SAHR）與 $\mathcal{M}_{\mathrm{off}\to\mathrm{relay}}(t)$（BPLR）。

\begin{figure*}[t]
  \centering
  \includegraphics[width=\textwidth]{ch4_beam_pair_mapping}
  \caption{Time slot $t$ 之 beam-pair mapping 對照表。(a) $\mathcal{M}_{\mathrm{safe}\to\mathrm{rel}}(t)$：$s^{\mathrm{c}}$ 上 safe beams 對應至 $s^{\mathrm{N}}$/$s^{\mathrm{S}}$ 之 helper release beams（SAHR）。(b) $\mathcal{M}_{\mathrm{off}\to\mathrm{relay}}(t)$：closed beams 對應至 helper relay beams（BPLR）。}
  \label{fig:ch4_beam_pair_mapping}
\end{figure*}

本文考慮的 Walker-star constellation 與 fixed beam layout 使 critical satellite 與 north/south helper satellites 呈現近似 one-to-one beam overlap。具體而言，critical satellite 的 north-half beams 主要與 north helper satellite 的 south-half beams 重疊，而 critical satellite 的 south-half beams 主要與 south helper satellite 的 north-half beams 重疊。因此，後續 relay 可在 beam pair 內局部執行，而不需進行 global matching。

此外，helper beams 分為 helper release beams 與 helper relay beams。前者用於釋放 helper satellite power，後者用於接手 critical closed-beam users。本文假設兩者於同一 time slot 中不重疊：
\begin{equation*}
\mathcal{B}_h^{\mathrm{rel}}(t)\cap \mathcal{B}_h^{\mathrm{relay}}(t)=\emptyset.
\end{equation*}

\subsection{Safe-Beam Assisted Helper Relief}

讀取 $\mathcal{R}_B(t)$ 後，online 程序首先執行 SAHR。圖~\ref{fig:ch4_Safe-beam_assisted_power_release} 示意 SAHR 流程。此階段的目的不是直接接手 critical closed-beam users，而是利用 critical safe beams 接手部分 helper users，使 helper satellite 釋放功率，供後續 relay boosting 使用。

對於
\begin{equation*}
(b_c^{\mathrm{safe}},b_h^{\mathrm{rel}})\in \mathcal{M}_{\mathrm{safe}\to\mathrm{rel}}(t),
\end{equation*}
系統找出目前由 $b_h^{\mathrm{rel}}$ 服務，且同時位於 $b_c^{\mathrm{safe}}$ footprint 內的 helper users：
\begin{equation*}
\begin{split}
&\mathcal{V}_{\mathrm{rel}}(b_c^{\mathrm{safe}},b_h^{\mathrm{rel}},t)= \\
&\quad
\left\{
v :
v\in\mathcal{U}_{s_h,b_h^{\mathrm{rel}}}(t),\;
v\in\mathcal{F}_{s^{\mathrm{c}},b_c^{\mathrm{safe}}}(t)
\right\}.
\end{split}
\end{equation*}
若將 candidate helper user $v$ reassociate 至 $b_c^{\mathrm{safe}}$，helper release beam 所需功率的下降量為
\begin{equation*}
\Delta P_v^{\mathrm{rel}}(t)=
P_{b_h^{\mathrm{rel}}}^{\mathrm{req,old}}(t)-
P_{b_h^{\mathrm{rel}}}^{\mathrm{req,new}}(v,t).
\end{equation*}
SAHR score 定義為
\begin{equation*}
S_{\mathrm{relief}}(v,t)=\Delta P_v^{\mathrm{rel}}(t).
\end{equation*}
candidate helper users 依據 $S_{\mathrm{relief}}(v,t)$ 由大到小排序。僅當被轉移 user 的 satisfaction 不下降，且 safe beam 原本 users 亦不下降時，才接受 reassociation：
\begin{equation*}
\mathrm{sat}_v^{\mathrm{safe}}(t)\ge \mathrm{sat}_v^{\mathrm{old}}(t),
\end{equation*}
\begin{equation*}
\mathrm{sat}_w^{\mathrm{new}}(t)\ge \mathrm{sat}_w^{\mathrm{old}}(t),
\qquad \forall w\in\mathcal{U}_{b_c^{\mathrm{safe}}}^{\mathrm{ori}}(t).
\end{equation*}
若條件成立，user $v$ 被 reassociate 至 critical safe beam，且 $\Delta P_v^{\mathrm{rel}}(t)$ 回到 helper satellite power pool。對於 helper satellite $s_h$，SAHR 釋放的總功率為
\begin{equation*}
P_{s_h}^{\mathrm{rel}}(t)=
\sum_{v\in\mathcal{A}_{\mathrm{rel}}(s_h,t)} \Delta P_v^{\mathrm{rel}}(t),
\end{equation*}
其中 $\mathcal{A}_{\mathrm{rel}}(s_h,t)$ 為成功被 safe beams 接手的 helper users 集合。

baseline operation 下原本未使用的功率為
\begin{equation*}
P_{s_h}^{\mathrm{unused}}(t)=
P_{s_h}^{\max}-\sum_{b\in\mathcal{B}_{s_h}} P^{\mathrm{base}},
\end{equation*}
因此，後續可用於 relay boosting 的 power pool 為
\begin{equation*}
P_{s_h}^{\mathrm{pool}}(t)=
P_{s_h}^{\mathrm{unused}}(t)+P_{s_h}^{\mathrm{rel}}(t).
\end{equation*}
除 helper release beams 與 helper relay beams 外，其餘 helper beams 皆維持 baseline power $P^{\mathrm{base}}$。

\begin{figure*}[t]
  \centering
  \includegraphics[width=\textwidth]{ch4_Safe-beam_assisted_power_release}
  \caption{Safe-Beam Assisted Helper Relief (SAHR) 示意。critical safe beams 接手 helper release beams 上符合條件的 users，所釋放功率回收到 helper satellite 的 power pool。}
  \label{fig:ch4_Safe-beam_assisted_power_release}
\end{figure*}

\subsection{Helper Relay Beam Boosting}

完成 SAHR 後，helper satellite $s_h$ 使用 $P_{s_h}^{\mathrm{pool}}(t)$ 提升 helper relay beams 的 transmit power。helper relay beams 由 $\mathcal{M}_{\mathrm{off}\to\mathrm{relay}}(t)$ 決定：
\begin{equation*}
\begin{split}
\mathcal{B}_{s_h}^{\mathrm{relay}}(t)=
\bigl\{
&b_h^{\mathrm{relay}} :
\exists\, b_c^{\mathrm{off}}\in\mathcal{B}_c^{\mathrm{off}}(t): \\
&\qquad (b_c^{\mathrm{off}},b_h^{\mathrm{relay}})\in \mathcal{M}_{\mathrm{off}\to\mathrm{relay}}(t)
\bigr\}.
\end{split}
\end{equation*}
將 relay beam 從 baseline power 提升至 beam upper bound 所需的額外功率為
\begin{equation*}
\begin{split}
\Delta P_{b_h^{\mathrm{relay}}}^{\max}(t)=
P_{s_h,b_h^{\mathrm{relay}}}^{\max}-P^{\mathrm{base}}.
\end{split}
\end{equation*}
對於每個 helper relay beam，其 relay demand 定義為
\begin{equation*}
D_{b_h^{\mathrm{relay}}}^{\mathrm{relay}}(t)=
\left|\mathcal{U}_{\mathrm{cand}}(b_h^{\mathrm{relay}},t)\right|,
\end{equation*}
其中 $\mathcal{U}_{\mathrm{cand}}(b_h^{\mathrm{relay}},t)$ 表示根據 $\mathcal{M}_{\mathrm{off}\to\mathrm{relay}}(t)$ 得到的 candidate critical users，也就是位於對應 critical closed beam 內，且可被 $b_h^{\mathrm{relay}}$ 覆蓋的 users。

若 $P_{s_h}^{\mathrm{pool}}(t)$ 足以支援所有 relay beams 提升至 upper bound，則所有 relay beams 皆可被 boost。若 power pool 不足，本文採用 relay-demand-priority boosting：helper relay beams 依據 $D_{b_h^{\mathrm{relay}}}^{\mathrm{relay}}(t)$ 由大到小排序，relay demand 較高者優先取得額外功率。

在分配過程中，若剩餘 power pool 足以使當前 relay beam 提升至 beam upper bound $P_{s_h,b_h^{\mathrm{relay}}}^{\max}$，則將其提升至上限。若剩餘 power pool 不足以完整提升下一個 relay beam，則仍可將剩餘功率分配給該 beam，使其 transmit power 高於 $P^{\mathrm{base}}$。如此可避免剩餘 power pool 被浪費，並提升該 relay beam 的可用 relay capacity。

因此，所有 helper relay beams 仍可進入後續 BPLR 階段，差別在於各 relay beams 依據最終 transmit power 具有不同的 available relay capacity。

\subsection{Beam-Pair Local Relay}

完成 relay-beam boosting 後，系統於每一組 closed-beam-to-helper-relay-beam pair 內執行 BPLR。圖~\ref{fig:ch4_Beam-pair_local_relay} 示意 BPLR 流程。對於
\begin{equation*}
(b_c^{\mathrm{off}},b_h^{\mathrm{relay}})\in \mathcal{M}_{\mathrm{off}\to\mathrm{relay}}(t),
\end{equation*}
系統找出原本由 $b_c^{\mathrm{off}}$ 服務，且同時位於 $b_h^{\mathrm{relay}}$ footprint 內的 critical users：
\begin{equation*}
\begin{split}
&\mathcal{U}_{\mathrm{pair}}(b_c^{\mathrm{off}},b_h^{\mathrm{relay}},t)= \\
&\quad
\left\{
u :
u\in\mathcal{U}_{s^{\mathrm{c}},b_c^{\mathrm{off}}}(t),\;
u\in\mathcal{F}_{s_h,b_h^{\mathrm{relay}}}(t)
\right\}.
\end{split}
\end{equation*}
對於 candidate critical user $u$，其由 helper relay beam 服務時的預期 achievable rate 記為
\begin{equation*}
R_{u,b_h^{\mathrm{relay}}}(t).
\end{equation*}
本文定義 BPLR score 為
\begin{equation*}
S_{\mathrm{relay}}(u,t)=R_{u,b_h^{\mathrm{relay}}}(t).
\end{equation*}
candidate users 依據 $S_{\mathrm{relay}}(u,t)$ 由大到小排序。排序後，helper relay beam 依據其 final transmit power 計算可用於 critical users 的 remaining relay capacity。此 capacity 已扣除保護原本 helper users 所需的服務資源。

接著，系統依序將 remaining relay capacity 分配給 sorted critical users。若剩餘容量足以滿足 user demand $D_u(t)$，則將該 user reassociate 至 helper relay beam，並更新其 satisfaction。若最後剩餘容量不足以完全滿足下一個 user，但仍大於零，則將剩餘容量分配給該 user，使其獲得 partial satisfaction。當 remaining relay capacity 用盡後，該 beam pair 的 relay process 即停止。

所有 relay decisions 仍需滿足 single association constraint。也就是說，被 relay 的 user 會從 critical closed beam reassociate 至 helper relay beam，且同一 time slot 中仍只由一個 beam 服務。

\begin{figure*}[t]
  \centering
  \includegraphics[width=\textwidth]{ch4_Beam-pair_local_relay}
  \caption{Beam-Pair Local Relay (BPLR) 示意。依 $\mathcal{M}_{\mathrm{off}\to\mathrm{relay}}(t)$ 將 closed beams 上的 candidate critical users reassociate 至對應的 helper relay beams。}
  \label{fig:ch4_Beam-pair_local_relay}
\end{figure*}

\subsection{Overall Algorithm}

本文之 online procedure 採用單一 time-slot-based algorithm。在每個 time slot 中，SAPR-R 先執行 SAHR 與 relay-beam boosting，接著根據更新後的 helper relay beam power 執行 BPLR。完整流程如 Algorithm~\ref{alg:ch4_sapr} 所示。

\begin{algorithm}[H]
\caption{SAPR-R Online Operation}
\label{alg:ch4_sapr}
\scriptsize
\begin{algorithmic}[1]
\Require $\{\mathcal{R}_B(t)\}_{t=1}^{T}$, initial $\mathbf{x}(t)$, $\mathcal{U}(t)$, $D_u(t)$
\Require $P^{\mathrm{base}}$, $P_{s,b}^{\max}$, $P_s^{\max}$, coverage, capacity/satisfaction model
\Ensure $\mathbf{x}^{\mathrm{final}}(t)$, $P_{s,b}^{\mathrm{final}}(t)$, $\mathrm{sat}_u^{\mathrm{final}}(t)$

\For{$t=1$ \textbf{to} $T$}
    \State Load $\mathcal{R}_B(t)$; initialize beam powers per record and baseline operation
    \State \textbf{Stage~1: SAHR and relay-beam boosting}
    \ForAll{$(b_c^{\mathrm{safe}},b_h^{\mathrm{rel}})\in \mathcal{M}_{\mathrm{safe}\to\mathrm{rel}}(t)$}
        \State Build candidate helper users; compute $\Delta P_v^{\mathrm{rel}}(t)$
        \State Sort candidates by $\Delta P_v^{\mathrm{rel}}(t)$ descending
        \ForAll{$v$ in sorted order}
            \State Trial reassociation $v$: $b_h^{\mathrm{rel}}\to b_c^{\mathrm{safe}}$
            \If{$v$ not degraded and users on $b_c^{\mathrm{safe}}$ not degraded}
                \State Accept; update association and $P_{s_h}^{\mathrm{pool}}(t)$
            \EndIf
        \EndFor
    \EndFor
    \ForAll{$s_h\in\{s^{\mathrm{N}}(t),s^{\mathrm{S}}(t)\}$}
        \State Extract $\mathcal{B}_{s_h}^{\mathrm{relay}}(t)$; compute relay demand; sort by demand
        \ForAll{$b_h^{\mathrm{relay}}\in \mathcal{B}_{s_h}^{\mathrm{relay}}(t)$ in sorted order}
            \State Allocate $P_{s_h}^{\mathrm{pool}}(t)$ (full boost, partial boost, or baseline)
        \EndFor
        \State Check $\sum_b P_{s,b}(t)\le P_{s_h}^{\max}$
    \EndFor
    \State \textbf{Stage~2: Beam-pair local relay}
    \ForAll{$(b_c^{\mathrm{off}},b_h^{\mathrm{relay}})\in \mathcal{M}_{\mathrm{off}\to\mathrm{relay}}(t)$}
        \State Build candidate critical users; sort by relay link quality
        \State Compute remaining relay capacity of $b_h^{\mathrm{relay}}$ (protect helper users)
        \ForAll{$u$ in sorted order}
            \If{capacity $\ge D_u(t)$}
                \State Reassociate $u$ to $b_h^{\mathrm{relay}}$; update $\mathrm{sat}_u(t)$ and capacity
            \ElsIf{remaining capacity $>0$}
                \State Partial allocation to $u$; set remaining capacity to $0$; \textbf{break}
            \Else
                \State \textbf{break}
            \EndIf
        \EndFor
    \EndFor
    \State Validate single association, beam power, satellite total power, and EPFD constraints
    \State Store $\mathbf{x}^{\mathrm{final}}(t)$, $P_{s,b}^{\mathrm{final}}(t)$, $\mathrm{sat}_u^{\mathrm{final}}(t)$
\EndFor
\end{algorithmic}
\end{algorithm}

\subsection{Complexity Discussion}

本文方法避免直接求解 mixed-integer nonlinear optimization problem。由於每個 time slot 的 beam relation record 已預先建立，online procedure 不需重新搜尋 critical satellite、helper satellites 或 beam-pair mappings，而是直接在對應 beam pairs 內執行局部 reassociation 與 relay decision。

令 $K_t^{\mathrm{rel}}$ 表示 time slot $t$ 中 SAHR 階段的 candidate helper users 數量，$K_t^{\mathrm{relay}}$ 表示 BPLR 階段的 candidate critical users 數量。在 Algorithm~\ref{alg:ch4_sapr} 中，candidate 建立與 released power 計算皆與 $K_t^{\mathrm{rel}}$ 呈線性關係，而主要成本來自根據 $\Delta P_v^{\mathrm{rel}}(t)$ 對 candidate helper users 排序。因此，SAHR 的時間複雜度為
\begin{equation*}
\mathcal{O}\!\bigl(K_t^{\mathrm{rel}}\log K_t^{\mathrm{rel}}\bigr).
\end{equation*}
在 BPLR 階段，candidate critical users 依據 relay link quality 排序，後續 capacity allocation 則為 sequential allocation。因此，BPLR 的時間複雜度為
\begin{equation*}
\mathcal{O}\!\bigl(K_t^{\mathrm{relay}}\log K_t^{\mathrm{relay}}\bigr).
\end{equation*}
綜合所有 time slots，本文方法的 online complexity 可表示為
\begin{equation*}
\mathcal{O}\!\left(
\sum_{t=1}^{T}
K_t^{\mathrm{rel}}\log K_t^{\mathrm{rel}}
+
K_t^{\mathrm{relay}}\log K_t^{\mathrm{relay}}
\right).
\end{equation*}
此結果顯示，本文方法將 online decision 限制於 precomputed beam pairs 內的 local sorting 與 sequential allocation，而非每個 time slot 重新執行 global user association 與 continuous power optimization。

\subsection{Summary}

本章提出 Safe-Beam Assisted Power-Release Relay (SAPR-R) 方法，以處理 critical satellite 接近 protected GS 時 EPFD constraint 與 user satisfaction 之間的 tradeoff。本文先利用 predictable satellite geometry 與 beam footprint overlap 建立 time-slot beam relation record，再於 online 階段執行 SAHR、relay beam boosting 與 BPLR。

透過此分階段設計，本文將原本高度耦合的 user association 與 beam power allocation 問題拆解為可局部處理的 decision steps。critical safe beams 用於釋放 helper satellite power，helper relay beams 則根據 relay demand 取得額外 transmit power，並於 beam pair 內接手 critical closed-beam users。此方法可在滿足 EPFD constraint 的前提下，提高 critical period 中的 total user satisfaction。
; compile chapter4_standalone.tex for preview.

% Reference style (like IEEE paper): **Label:** regular body text on the same line.
% - Label before colon: bold (bfseries)
% - Following text: normal weight (mdseries / normalfont)
\makeatletter
\@ifundefined{md@runinlabel}{%
  \newcommand{\md@runinlabel}[1]{{\bfseries #1}\normalfont}%
  \newcommand{\mdsubsubsection}[2][]{%
    \refstepcounter{subsubsection}%
    \par\vspace{0.35ex}%
    \noindent\md@runinlabel{\arabic{subsubsection})~#2:}\space
    \def\@mdss@label{#1}%
    \ifx\@mdss@label\@empty\else\label{#1}\fi
  }%
  \newcommand{\mdparagraph}[2][]{%
    \par\vspace{0.25ex}%
    \noindent\md@runinlabel{#2:}\space
    \def\@mdpar@label{#1}%
    \ifx\@mdpar@label\@empty\else\label{#1}\fi
  }%
}{}
\makeatother


\section{Method}
\label{sec:chapter4}

\subsection{Method Overview}

第三章已將本文問題表述為在 user association、beam power、satellite total power 與 EPFD constraints 下最大化 total user satisfaction 的最佳化問題。然而，該問題同時包含離散的 user--beam association decision 與連續的 beam power configuration，且 association、SINR、TDMA sharing、satisfaction 與 aggregate EPFD 彼此耦合，因此直接求解全域最佳解具有較高計算複雜度。

為降低 online decision complexity，本文提出一套 Safe-Beam Assisted Power-Release Relay (SAPR-R) 方法。SAPR-R 利用 LEO constellation、protected GS location、GSO direction 與 fixed beam layout 的可預測性，先針對整個觀察期間的所有 time slots 建立 beam-level relation records。對於每個 time slot，系統預先記錄 critical satellite、north/south helper satellites、critical satellite 上需關閉的 beams、仍可運作的 safe beams，以及 critical beams 與 helper beams 之間的 beam-pair mappings。此階段僅依賴 satellite geometry、beam footprint 與 EPFD calculation，不需要 user distribution，因此可於 online user reassociation 前完成。

在 online operation 階段，系統讀取當前 time slot 的 beam relation record，並根據當下 users 的位置與 association 執行 user-aware relay decision。首先，SAPR-R 透過 Safe-Beam Assisted Helper Relief (SAHR) 使 critical safe beams 嘗試接手 helper release beams 上的部分 users。若 reassociation 不降低被轉移 user 與 safe beam 原本 users 的 satisfaction，則接受此次轉移，並將 helper release beam 所釋放出的功率回收到 helper satellite 的 available power pool。接著，helper satellite 依據 relay demand 將 power pool 分配給 helper relay beams，以提升其後續接手 closed-beam users 的 capacity。最後，SAPR-R 透過 Beam-Pair Local Relay (BPLR)，在每組 closed beam/helper relay beam pair 內，依據 relay link quality 對 critical closed-beam users 進行排序與 sequential allocation。

綜合而言，本文方法包含下列步驟：
\begin{itemize}[leftmargin=*,nosep]
  \item \textbf{Time-slot beam relation record:} 預先記錄每個 time slot 的 critical/helper satellites、closed/safe beams 與 beam-pair mappings。
  \item \textbf{Safe-Beam Assisted Helper Relief (SAHR):} 利用 critical safe beams 接手部分 helper users，在不降低 satisfaction 的條件下釋放 helper satellite power。
  \item \textbf{Helper relay beam boosting:} 依據 relay demand 將 available power pool 分配給 helper relay beams。
  \item \textbf{Beam-Pair Local Relay (BPLR):} 在每組 beam pair 內，依 relay link quality 分配 relay capacity 給 closed-beam users。
\end{itemize}

\subsection{Time-Slot Beam Relation Record}

\begin{figure*}[t]
  \centering
  \includegraphics[width=\textwidth]{ch4_beam_relation_record}
  \caption{三個 time slots 中，沿軌道方向之三顆 LEO 衛星（$s^{\mathrm{N}}$、$s^{\mathrm{c}}$、$s^{\mathrm{S}}$）飛過受保護 GSO 地面站之示意。critical 衛星區域 $S_{\mathrm{C}}$ 內，黑色斜線為 $\mathcal{B}_c^{\mathrm{off}}(t)$（closed beams），綠色斜線為 $\mathcal{B}_c^{\mathrm{safe}}(t)$（safe beams）。(a)--(c) 為由早至晚之 time slot。}
  \label{fig:ch4_beam_relation_record}
\end{figure*}

對於每一個 time slot $t$，本文建立一個 beam relation record $\mathcal{R}_B(t)$。圖~\ref{fig:ch4_beam_relation_record} 示意本文所考慮之 time-slot beam 狀態，包括：
\begin{equation*}
\begin{split}
\mathcal{R}_B(t)=\bigl(
& s^{\mathrm{c}}(t),\,
s^{\mathrm{N}}(t),\,
s^{\mathrm{S}}(t), \\
& \mathcal{B}_c^{\mathrm{off}}(t),\,
\mathcal{B}_c^{\mathrm{safe}}(t), \\
& \mathcal{M}_{\mathrm{safe}\to\mathrm{rel}}(t),\,
\mathcal{M}_{\mathrm{off}\to\mathrm{relay}}(t)
\bigr).
\end{split}
\end{equation*}
該 record 儲存後續 SAHR 與 BPLR 所需的 beam-level information，其中 $s^{\mathrm{c}}(t)$、$s^{\mathrm{N}}(t)$ 與 $s^{\mathrm{S}}(t)$ 分別為 critical satellite、north helper satellite 與 south helper satellite。$\mathcal{B}_c^{\mathrm{off}}(t)$ 為 critical satellite 上因 EPFD constraint 需被關閉的 beams，$\mathcal{B}_c^{\mathrm{safe}}(t)$ 為未被關閉且維持 baseline power 的 safe beams。$\mathcal{M}_{\mathrm{safe}\to\mathrm{rel}}(t)$ 與 $\mathcal{M}_{\mathrm{off}\to\mathrm{relay}}(t)$ 分別表示 safe-beam-to-helper-release-beam mapping 與 closed-beam-to-helper-relay-beam mapping。

\mdsubsubsection{Critical and Helper Satellite Identification}
在每個 time slot $t$，系統根據 LEO satellite 與 protected GS 之間的幾何關係，找出對 GS 造成主要 EPFD 風險的 critical satellite：
\begin{equation*}
s^{\mathrm{c}}(t).
\end{equation*}
根據模擬觀察，EPFD violation 主要由 critical satellite 上的 high-risk beams 主導，因此本文僅針對 $s^{\mathrm{c}}(t)$ 上的 beams 進行 shutdown decision，不關閉 helper satellites 或其他 visible satellites 的 beams。接著，選取 critical satellite 同軌道方向上的 north neighbor 與 south neighbor 作為 helper satellites：
\begin{equation*}
s^{\mathrm{N}}(t),\quad s^{\mathrm{S}}(t).
\end{equation*}

\mdsubsubsection{EPFD-Driven Beam Status Determination}
對每個 critical beam $b\in\mathcal{B}_{s^{\mathrm{c}}}$，定義其 EPFD risk score 為
\begin{equation*}
q_{c,b}(t)=\mathrm{EPFD}_{s^{\mathrm{c}},b}(t).
\end{equation*}
critical beams 依據 $q_{c,b}(t)$ 由大到小排序，並採用 greedy shutdown rule：每次關閉目前 EPFD risk score 最高的 beam，直到 aggregate EPFD 滿足
\begin{equation*}
\mathrm{EPFD}_{\mathrm{agg}}(t)\le \mathrm{EPFD}_{\mathrm{lim}}.
\end{equation*}
被關閉的 beams 組成
\begin{equation*}
\mathcal{B}_c^{\mathrm{off}}(t),
\end{equation*}
其餘 beams 則形成 safe beam set：
\begin{equation*}
\mathcal{B}_c^{\mathrm{safe}}(t)=\mathcal{B}_{s^{\mathrm{c}}}\setminus \mathcal{B}_c^{\mathrm{off}}(t).
\end{equation*}
此階段僅根據 beam-level EPFD contribution 與 regulatory limit 決定 beam status，不涉及 user distribution，因此可於 online user reassociation 前預先完成。

\mdsubsubsection{Beam-Pair Mapping Construction}
在得到 $\mathcal{B}_c^{\mathrm{off}}(t)$ 與 $\mathcal{B}_c^{\mathrm{safe}}(t)$ 後，本文根據 beam footprint overlap 建立 beam-pair mappings。對於 SAHR，定義
\begin{equation*}
\begin{split}
\mathcal{M}_{\mathrm{safe}\to\mathrm{rel}}(t)=
\bigl\{
&(b_c^{\mathrm{safe}},b_h^{\mathrm{rel}}) : \\
&\mathcal{F}_{s^{\mathrm{c}},b_c^{\mathrm{safe}}}(t)\cap
\mathcal{F}_{s_h,b_h^{\mathrm{rel}}}(t)\neq\emptyset, \\
& s_h\in\{s^{\mathrm{N}}(t),s^{\mathrm{S}}(t)\}
\bigr\},
\end{split}
\end{equation*}
其中 $s_h\in\{s^{\mathrm{N}}(t),s^{\mathrm{S}}(t)\}$。此 mapping 用於找出 critical safe beam 可接手哪些 helper release beam 上的 users。

對於 BPLR，定義
\begin{equation*}
\begin{split}
\mathcal{M}_{\mathrm{off}\to\mathrm{relay}}(t)=
\bigl\{
&(b_c^{\mathrm{off}},b_h^{\mathrm{relay}}) : \\
&\mathcal{F}_{s^{\mathrm{c}},b_c^{\mathrm{off}}}(t)\cap
\mathcal{F}_{s_h,b_h^{\mathrm{relay}}}(t)\neq\emptyset
\bigr\}.
\end{split}
\end{equation*}
此 mapping 用於找出 critical closed-beam users 可由哪些 helper relay beams 接手。

圖~\ref{fig:ch4_beam_pair_mapping} 呈現儲存於 $\mathcal{R}_B(t)$ 中的兩組 beam-pair mappings：$\mathcal{M}_{\mathrm{safe}\to\mathrm{rel}}(t)$（SAHR）與 $\mathcal{M}_{\mathrm{off}\to\mathrm{relay}}(t)$（BPLR）。

\begin{figure*}[t]
  \centering
  \includegraphics[width=\textwidth]{ch4_beam_pair_mapping}
  \caption{Time slot $t$ 之 beam-pair mapping 對照表。(a) $\mathcal{M}_{\mathrm{safe}\to\mathrm{rel}}(t)$：$s^{\mathrm{c}}$ 上 safe beams 對應至 $s^{\mathrm{N}}$/$s^{\mathrm{S}}$ 之 helper release beams（SAHR）。(b) $\mathcal{M}_{\mathrm{off}\to\mathrm{relay}}(t)$：closed beams 對應至 helper relay beams（BPLR）。}
  \label{fig:ch4_beam_pair_mapping}
\end{figure*}

本文考慮的 Walker-star constellation 與 fixed beam layout 使 critical satellite 與 north/south helper satellites 呈現近似 one-to-one beam overlap。具體而言，critical satellite 的 north-half beams 主要與 north helper satellite 的 south-half beams 重疊，而 critical satellite 的 south-half beams 主要與 south helper satellite 的 north-half beams 重疊。因此，後續 relay 可在 beam pair 內局部執行，而不需進行 global matching。

此外，helper beams 分為 helper release beams 與 helper relay beams。前者用於釋放 helper satellite power，後者用於接手 critical closed-beam users。本文假設兩者於同一 time slot 中不重疊：
\begin{equation*}
\mathcal{B}_h^{\mathrm{rel}}(t)\cap \mathcal{B}_h^{\mathrm{relay}}(t)=\emptyset.
\end{equation*}

\subsection{Safe-Beam Assisted Helper Relief}

讀取 $\mathcal{R}_B(t)$ 後，online 程序首先執行 SAHR。圖~\ref{fig:ch4_Safe-beam_assisted_power_release} 示意 SAHR 流程。此階段的目的不是直接接手 critical closed-beam users，而是利用 critical safe beams 接手部分 helper users，使 helper satellite 釋放功率，供後續 relay boosting 使用。

對於
\begin{equation*}
(b_c^{\mathrm{safe}},b_h^{\mathrm{rel}})\in \mathcal{M}_{\mathrm{safe}\to\mathrm{rel}}(t),
\end{equation*}
系統找出目前由 $b_h^{\mathrm{rel}}$ 服務，且同時位於 $b_c^{\mathrm{safe}}$ footprint 內的 helper users：
\begin{equation*}
\begin{split}
&\mathcal{V}_{\mathrm{rel}}(b_c^{\mathrm{safe}},b_h^{\mathrm{rel}},t)= \\
&\quad
\left\{
v :
v\in\mathcal{U}_{s_h,b_h^{\mathrm{rel}}}(t),\;
v\in\mathcal{F}_{s^{\mathrm{c}},b_c^{\mathrm{safe}}}(t)
\right\}.
\end{split}
\end{equation*}
若將 candidate helper user $v$ reassociate 至 $b_c^{\mathrm{safe}}$，helper release beam 所需功率的下降量為
\begin{equation*}
\Delta P_v^{\mathrm{rel}}(t)=
P_{b_h^{\mathrm{rel}}}^{\mathrm{req,old}}(t)-
P_{b_h^{\mathrm{rel}}}^{\mathrm{req,new}}(v,t).
\end{equation*}
SAHR score 定義為
\begin{equation*}
S_{\mathrm{relief}}(v,t)=\Delta P_v^{\mathrm{rel}}(t).
\end{equation*}
candidate helper users 依據 $S_{\mathrm{relief}}(v,t)$ 由大到小排序。僅當被轉移 user 的 satisfaction 不下降，且 safe beam 原本 users 亦不下降時，才接受 reassociation：
\begin{equation*}
\mathrm{sat}_v^{\mathrm{safe}}(t)\ge \mathrm{sat}_v^{\mathrm{old}}(t),
\end{equation*}
\begin{equation*}
\mathrm{sat}_w^{\mathrm{new}}(t)\ge \mathrm{sat}_w^{\mathrm{old}}(t),
\qquad \forall w\in\mathcal{U}_{b_c^{\mathrm{safe}}}^{\mathrm{ori}}(t).
\end{equation*}
若條件成立，user $v$ 被 reassociate 至 critical safe beam，且 $\Delta P_v^{\mathrm{rel}}(t)$ 回到 helper satellite power pool。對於 helper satellite $s_h$，SAHR 釋放的總功率為
\begin{equation*}
P_{s_h}^{\mathrm{rel}}(t)=
\sum_{v\in\mathcal{A}_{\mathrm{rel}}(s_h,t)} \Delta P_v^{\mathrm{rel}}(t),
\end{equation*}
其中 $\mathcal{A}_{\mathrm{rel}}(s_h,t)$ 為成功被 safe beams 接手的 helper users 集合。

baseline operation 下原本未使用的功率為
\begin{equation*}
P_{s_h}^{\mathrm{unused}}(t)=
P_{s_h}^{\max}-\sum_{b\in\mathcal{B}_{s_h}} P^{\mathrm{base}},
\end{equation*}
因此，後續可用於 relay boosting 的 power pool 為
\begin{equation*}
P_{s_h}^{\mathrm{pool}}(t)=
P_{s_h}^{\mathrm{unused}}(t)+P_{s_h}^{\mathrm{rel}}(t).
\end{equation*}
除 helper release beams 與 helper relay beams 外，其餘 helper beams 皆維持 baseline power $P^{\mathrm{base}}$。

\begin{figure*}[t]
  \centering
  \includegraphics[width=\textwidth]{ch4_Safe-beam_assisted_power_release}
  \caption{Safe-Beam Assisted Helper Relief (SAHR) 示意。critical safe beams 接手 helper release beams 上符合條件的 users，所釋放功率回收到 helper satellite 的 power pool。}
  \label{fig:ch4_Safe-beam_assisted_power_release}
\end{figure*}

\subsection{Helper Relay Beam Boosting}

完成 SAHR 後，helper satellite $s_h$ 使用 $P_{s_h}^{\mathrm{pool}}(t)$ 提升 helper relay beams 的 transmit power。helper relay beams 由 $\mathcal{M}_{\mathrm{off}\to\mathrm{relay}}(t)$ 決定：
\begin{equation*}
\begin{split}
\mathcal{B}_{s_h}^{\mathrm{relay}}(t)=
\bigl\{
&b_h^{\mathrm{relay}} :
\exists\, b_c^{\mathrm{off}}\in\mathcal{B}_c^{\mathrm{off}}(t): \\
&\qquad (b_c^{\mathrm{off}},b_h^{\mathrm{relay}})\in \mathcal{M}_{\mathrm{off}\to\mathrm{relay}}(t)
\bigr\}.
\end{split}
\end{equation*}
將 relay beam 從 baseline power 提升至 beam upper bound 所需的額外功率為
\begin{equation*}
\begin{split}
\Delta P_{b_h^{\mathrm{relay}}}^{\max}(t)=
P_{s_h,b_h^{\mathrm{relay}}}^{\max}-P^{\mathrm{base}}.
\end{split}
\end{equation*}
對於每個 helper relay beam，其 relay demand 定義為
\begin{equation*}
D_{b_h^{\mathrm{relay}}}^{\mathrm{relay}}(t)=
\left|\mathcal{U}_{\mathrm{cand}}(b_h^{\mathrm{relay}},t)\right|,
\end{equation*}
其中 $\mathcal{U}_{\mathrm{cand}}(b_h^{\mathrm{relay}},t)$ 表示根據 $\mathcal{M}_{\mathrm{off}\to\mathrm{relay}}(t)$ 得到的 candidate critical users，也就是位於對應 critical closed beam 內，且可被 $b_h^{\mathrm{relay}}$ 覆蓋的 users。

若 $P_{s_h}^{\mathrm{pool}}(t)$ 足以支援所有 relay beams 提升至 upper bound，則所有 relay beams 皆可被 boost。若 power pool 不足，本文採用 relay-demand-priority boosting：helper relay beams 依據 $D_{b_h^{\mathrm{relay}}}^{\mathrm{relay}}(t)$ 由大到小排序，relay demand 較高者優先取得額外功率。

在分配過程中，若剩餘 power pool 足以使當前 relay beam 提升至 beam upper bound $P_{s_h,b_h^{\mathrm{relay}}}^{\max}$，則將其提升至上限。若剩餘 power pool 不足以完整提升下一個 relay beam，則仍可將剩餘功率分配給該 beam，使其 transmit power 高於 $P^{\mathrm{base}}$。如此可避免剩餘 power pool 被浪費，並提升該 relay beam 的可用 relay capacity。

因此，所有 helper relay beams 仍可進入後續 BPLR 階段，差別在於各 relay beams 依據最終 transmit power 具有不同的 available relay capacity。

\subsection{Beam-Pair Local Relay}

完成 relay-beam boosting 後，系統於每一組 closed-beam-to-helper-relay-beam pair 內執行 BPLR。圖~\ref{fig:ch4_Beam-pair_local_relay} 示意 BPLR 流程。對於
\begin{equation*}
(b_c^{\mathrm{off}},b_h^{\mathrm{relay}})\in \mathcal{M}_{\mathrm{off}\to\mathrm{relay}}(t),
\end{equation*}
系統找出原本由 $b_c^{\mathrm{off}}$ 服務，且同時位於 $b_h^{\mathrm{relay}}$ footprint 內的 critical users：
\begin{equation*}
\begin{split}
&\mathcal{U}_{\mathrm{pair}}(b_c^{\mathrm{off}},b_h^{\mathrm{relay}},t)= \\
&\quad
\left\{
u :
u\in\mathcal{U}_{s^{\mathrm{c}},b_c^{\mathrm{off}}}(t),\;
u\in\mathcal{F}_{s_h,b_h^{\mathrm{relay}}}(t)
\right\}.
\end{split}
\end{equation*}
對於 candidate critical user $u$，其由 helper relay beam 服務時的預期 achievable rate 記為
\begin{equation*}
R_{u,b_h^{\mathrm{relay}}}(t).
\end{equation*}
本文定義 BPLR score 為
\begin{equation*}
S_{\mathrm{relay}}(u,t)=R_{u,b_h^{\mathrm{relay}}}(t).
\end{equation*}
candidate users 依據 $S_{\mathrm{relay}}(u,t)$ 由大到小排序。排序後，helper relay beam 依據其 final transmit power 計算可用於 critical users 的 remaining relay capacity。此 capacity 已扣除保護原本 helper users 所需的服務資源。

接著，系統依序將 remaining relay capacity 分配給 sorted critical users。若剩餘容量足以滿足 user demand $D_u(t)$，則將該 user reassociate 至 helper relay beam，並更新其 satisfaction。若最後剩餘容量不足以完全滿足下一個 user，但仍大於零，則將剩餘容量分配給該 user，使其獲得 partial satisfaction。當 remaining relay capacity 用盡後，該 beam pair 的 relay process 即停止。

所有 relay decisions 仍需滿足 single association constraint。也就是說，被 relay 的 user 會從 critical closed beam reassociate 至 helper relay beam，且同一 time slot 中仍只由一個 beam 服務。

\begin{figure*}[t]
  \centering
  \includegraphics[width=\textwidth]{ch4_Beam-pair_local_relay}
  \caption{Beam-Pair Local Relay (BPLR) 示意。依 $\mathcal{M}_{\mathrm{off}\to\mathrm{relay}}(t)$ 將 closed beams 上的 candidate critical users reassociate 至對應的 helper relay beams。}
  \label{fig:ch4_Beam-pair_local_relay}
\end{figure*}

\subsection{Overall Algorithm}

本文之 online procedure 採用單一 time-slot-based algorithm。在每個 time slot 中，SAPR-R 先執行 SAHR 與 relay-beam boosting，接著根據更新後的 helper relay beam power 執行 BPLR。完整流程如 Algorithm~\ref{alg:ch4_sapr} 所示。

\begin{algorithm}[H]
\caption{SAPR-R Online Operation}
\label{alg:ch4_sapr}
\scriptsize
\begin{algorithmic}[1]
\Require $\{\mathcal{R}_B(t)\}_{t=1}^{T}$, initial $\mathbf{x}(t)$, $\mathcal{U}(t)$, $D_u(t)$
\Require $P^{\mathrm{base}}$, $P_{s,b}^{\max}$, $P_s^{\max}$, coverage, capacity/satisfaction model
\Ensure $\mathbf{x}^{\mathrm{final}}(t)$, $P_{s,b}^{\mathrm{final}}(t)$, $\mathrm{sat}_u^{\mathrm{final}}(t)$

\For{$t=1$ \textbf{to} $T$}
    \State Load $\mathcal{R}_B(t)$; initialize beam powers per record and baseline operation
    \State \textbf{Stage~1: SAHR and relay-beam boosting}
    \ForAll{$(b_c^{\mathrm{safe}},b_h^{\mathrm{rel}})\in \mathcal{M}_{\mathrm{safe}\to\mathrm{rel}}(t)$}
        \State Build candidate helper users; compute $\Delta P_v^{\mathrm{rel}}(t)$
        \State Sort candidates by $\Delta P_v^{\mathrm{rel}}(t)$ descending
        \ForAll{$v$ in sorted order}
            \State Trial reassociation $v$: $b_h^{\mathrm{rel}}\to b_c^{\mathrm{safe}}$
            \If{$v$ not degraded and users on $b_c^{\mathrm{safe}}$ not degraded}
                \State Accept; update association and $P_{s_h}^{\mathrm{pool}}(t)$
            \EndIf
        \EndFor
    \EndFor
    \ForAll{$s_h\in\{s^{\mathrm{N}}(t),s^{\mathrm{S}}(t)\}$}
        \State Extract $\mathcal{B}_{s_h}^{\mathrm{relay}}(t)$; compute relay demand; sort by demand
        \ForAll{$b_h^{\mathrm{relay}}\in \mathcal{B}_{s_h}^{\mathrm{relay}}(t)$ in sorted order}
            \State Allocate $P_{s_h}^{\mathrm{pool}}(t)$ (full boost, partial boost, or baseline)
        \EndFor
        \State Check $\sum_b P_{s,b}(t)\le P_{s_h}^{\max}$
    \EndFor
    \State \textbf{Stage~2: Beam-pair local relay}
    \ForAll{$(b_c^{\mathrm{off}},b_h^{\mathrm{relay}})\in \mathcal{M}_{\mathrm{off}\to\mathrm{relay}}(t)$}
        \State Build candidate critical users; sort by relay link quality
        \State Compute remaining relay capacity of $b_h^{\mathrm{relay}}$ (protect helper users)
        \ForAll{$u$ in sorted order}
            \If{capacity $\ge D_u(t)$}
                \State Reassociate $u$ to $b_h^{\mathrm{relay}}$; update $\mathrm{sat}_u(t)$ and capacity
            \ElsIf{remaining capacity $>0$}
                \State Partial allocation to $u$; set remaining capacity to $0$; \textbf{break}
            \Else
                \State \textbf{break}
            \EndIf
        \EndFor
    \EndFor
    \State Validate single association, beam power, satellite total power, and EPFD constraints
    \State Store $\mathbf{x}^{\mathrm{final}}(t)$, $P_{s,b}^{\mathrm{final}}(t)$, $\mathrm{sat}_u^{\mathrm{final}}(t)$
\EndFor
\end{algorithmic}
\end{algorithm}

\subsection{Complexity Discussion}

本文方法避免直接求解 mixed-integer nonlinear optimization problem。由於每個 time slot 的 beam relation record 已預先建立，online procedure 不需重新搜尋 critical satellite、helper satellites 或 beam-pair mappings，而是直接在對應 beam pairs 內執行局部 reassociation 與 relay decision。

令 $K_t^{\mathrm{rel}}$ 表示 time slot $t$ 中 SAHR 階段的 candidate helper users 數量，$K_t^{\mathrm{relay}}$ 表示 BPLR 階段的 candidate critical users 數量。在 Algorithm~\ref{alg:ch4_sapr} 中，candidate 建立與 released power 計算皆與 $K_t^{\mathrm{rel}}$ 呈線性關係，而主要成本來自根據 $\Delta P_v^{\mathrm{rel}}(t)$ 對 candidate helper users 排序。因此，SAHR 的時間複雜度為
\begin{equation*}
\mathcal{O}\!\bigl(K_t^{\mathrm{rel}}\log K_t^{\mathrm{rel}}\bigr).
\end{equation*}
在 BPLR 階段，candidate critical users 依據 relay link quality 排序，後續 capacity allocation 則為 sequential allocation。因此，BPLR 的時間複雜度為
\begin{equation*}
\mathcal{O}\!\bigl(K_t^{\mathrm{relay}}\log K_t^{\mathrm{relay}}\bigr).
\end{equation*}
綜合所有 time slots，本文方法的 online complexity 可表示為
\begin{equation*}
\mathcal{O}\!\left(
\sum_{t=1}^{T}
K_t^{\mathrm{rel}}\log K_t^{\mathrm{rel}}
+
K_t^{\mathrm{relay}}\log K_t^{\mathrm{relay}}
\right).
\end{equation*}
此結果顯示，本文方法將 online decision 限制於 precomputed beam pairs 內的 local sorting 與 sequential allocation，而非每個 time slot 重新執行 global user association 與 continuous power optimization。

\subsection{Summary}

本章提出 Safe-Beam Assisted Power-Release Relay (SAPR-R) 方法，以處理 critical satellite 接近 protected GS 時 EPFD constraint 與 user satisfaction 之間的 tradeoff。本文先利用 predictable satellite geometry 與 beam footprint overlap 建立 time-slot beam relation record，再於 online 階段執行 SAHR、relay beam boosting 與 BPLR。

透過此分階段設計，本文將原本高度耦合的 user association 與 beam power allocation 問題拆解為可局部處理的 decision steps。critical safe beams 用於釋放 helper satellite power，helper relay beams 則根據 relay demand 取得額外 transmit power，並於 beam pair 內接手 critical closed-beam users。此方法可在滿足 EPFD constraint 的前提下，提高 critical period 中的 total user satisfaction。

\bibliographystyle{IEEEtran}
\bibliography{references}

\end{document}
